% =====================================================================
%  Statement-Design Arbitrage and Verification Limits:
%  A Forensic Code Audit of the openai/NavierStokesAndEuler Repository
%
%  Companion to ns-audit/AUDIT_LOG.md, entries E0--E148.
%  All code citations are to commit f9e8bc5.
% =====================================================================
% Compile with LuaLaTeX or XeLaTeX: the Lean listings are quoted verbatim,
% including Unicode connectives and blackboard-bold symbols.
\documentclass[11pt]{article}

\usepackage[margin=1in]{geometry}
% LuaLaTeX is required, not merely preferred: the Lean listings quote Unicode.
% The default Latin Modern mono has no ∀/∃/ℝ/ν/∧, which the first typeset run
% reported as 45 "Missing character" warnings and rendered as blanks. DejaVu Sans
% Mono covers them; FreeMono is the fallback if the image lacks DejaVu.
\usepackage{fontspec}
\IfFontExistsTF{DejaVu Sans Mono}
  {\setmonofont{DejaVu Sans Mono}[Scale=0.85]}
  {\IfFontExistsTF{FreeMono}{\setmonofont{FreeMono}[Scale=0.85]}{}}
\usepackage{amsmath,amssymb,amsthm,mathtools}
\usepackage{booktabs}
\usepackage{listings}
\usepackage{xcolor}
\usepackage{microtype}
\usepackage[colorlinks=true,linkcolor=blue!55!black,citecolor=blue!55!black,
            urlcolor=blue!55!black]{hyperref}

% ---------------------------------------------------------------- macros
\newcommand{\R}{\mathbb{R}}
\newcommand{\N}{\mathbb{N}}
\newcommand{\T}{\mathbb{T}}
\newcommand{\dd}{\,\mathrm{d}}
\newcommand{\norm}[1]{\left\lVert #1\right\rVert}
\newcommand{\abs}[1]{\left\lvert #1\right\rvert}
\newcommand{\supp}{\operatorname{tsupp}}
\newcommand{\Ph}{\Phi}
\newcommand{\ct}[1]{\texttt{#1}}
\newcommand{\Lean}{\textsc{Lean}}
\newcommand{\file}[1]{\texttt{\small #1}}
\DeclareMathOperator*{\esssup}{ess\,sup}
\DeclareMathOperator{\diver}{div}
\DeclareMathOperator{\loc}{loc}

\theoremstyle{plain}
\newtheorem{theorem}{Theorem}[section]
\newtheorem{lemma}[theorem]{Lemma}
\newtheorem{proposition}[theorem]{Proposition}
\newtheorem{corollary}[theorem]{Corollary}
\theoremstyle{definition}
\newtheorem{definition}[theorem]{Definition}
\newtheorem{example}[theorem]{Example}
\theoremstyle{remark}
\newtheorem{remark}[theorem]{Remark}
\newtheorem{claim}[theorem]{Claim}

% ------------------------------------------------------------- listings
\definecolor{leankw}{RGB}{20,80,150}
\definecolor{leancom}{RGB}{90,110,90}
\definecolor{leanbg}{RGB}{248,248,245}
\lstdefinelanguage{Lean4}{
  morekeywords={theorem,lemma,def,noncomputable,structure,namespace,end,open,
    import,section,by,exact,apply,intro,refine,obtain,have,let,show,calc,
    fun,forall,exists,Prop,Type,where,variable},
  sensitive=true,
  morecomment=[l]{--},
  morecomment=[s]{/-}{-/},
  morestring=[b]"
}
\lstset{
  language=Lean4,
  basicstyle=\ttfamily\footnotesize,
  keywordstyle=\color{leankw}\bfseries,
  commentstyle=\color{leancom}\itshape,
  backgroundcolor=\color{leanbg},
  frame=single, framesep=4pt, rulecolor=\color{black!25},
  breaklines=true, breakatwhitespace=true,
  columns=fullflexible, keepspaces=true,
  xleftmargin=8pt, xrightmargin=4pt,
  aboveskip=8pt, belowskip=8pt,
  literate={←}{{$\leftarrow$}}1 {→}{{$\rightarrow$}}1 {≤}{{$\leq$}}1
           {≥}{{$\geq$}}1 {≠}{{$\neq$}}1 {∀}{{$\forall$}}1 {∃}{{$\exists$}}1
           {∈}{{$\in$}}1 {∧}{{$\wedge$}}1 {∨}{{$\vee$}}1 {¬}{{$\neg$}}1
           {ℝ}{{$\mathbb{R}$}}2 {ℂ}{{$\mathbb{C}$}}2 {ℕ}{{$\mathbb{N}$}}2
           {𝓢}{{$\mathcal{S}$}}2 {ν}{{$\nu$}}1 {∞}{{$\infty$}}1 {ᵐ}{{$_{\mathrm{a.e.}}$}}2
           {⁻¹}{{$^{-1}$}}2 {·}{{$\cdot$}}1 {ᶠ}{{$_{\mathrm{f}}$}}2
}

\setlength{\emergencystretch}{2em}

\title{\bfseries Statement-Design Arbitrage and Verification Limits:\\
A Forensic Code Audit of the \texttt{openai/NavierStokesAndEuler} Repository}
\author{Abdelrahman El Sayad \and John}
\date{November 2026}

\begin{document}
\maketitle

\begin{abstract}
The repository \texttt{openai/NavierStokesAndEuler} (commit \texttt{f9e8bc5};
$641{,}332$ lines across $2{,}659$ files) contains a machine-checked proof of a
theorem asserting the existence, for every $\nu>0$, of a smooth divergence-free
velocity field with compact spatial support, zero initial datum, bounded kinetic
energy, and unbounded $L^\infty$ norm at $t=1$, together with the nonexistence
of any global finite-energy solution of the same forced problem. Macroscopic
caveats of that statement---the identically zero initial datum, and a forcing
term defined as the residual of the constructed field---have been noted in
public discussion and are consistent with the repository's own documentation;
we do not re-litigate them here.

This paper instead isolates two \emph{structural} mechanisms that explain how
proof closure is achieved, and what that closure does and does not certify.
The first is the \emph{Non-Evaluated Constant Barrier}: every quantitative
constant in the analytic chain is introduced by an existential binder or by
\ct{Classical.choose}, so the kernel verifies that some constant exists and
never that it is the right size or dimension. The barrier is a consequence of
which quantities are reflected into \emph{types}: exponents that appear as
indices (the $\ell^p$ exponent, the differentiability order) are checked by
unification and a wrong one is a type error, whereas a multiplicative constant
inhabiting $\exists C\ge 0,\ \forall x,\ \norm{Tx}\le C\norm{Sx}$ is checked
only for inhabitation. Unification decides agreement with the \emph{declared}
type family, not the fitness of that family as a model of the intended problem:
a consistently misspecified exponent elaborates without complaint, so the
guarantee is internal consistency, not correspondence. We show this is
epistemic, not logical: the headline
theorem mentions no constant, so the barrier cannot make it false, but it
forecloses effective bounds, machine dimension analysis, and any claim that the
constants are small enough for a subsequent numerical step.

The second is \emph{Anisotropic Viscous Scaling under the generalized cone}:
the repository's own pullback identities show that $\partial_s^2$ multiplies the
profile weight $q^{b}$ by $q^{-2}$ while $\partial_z^2$ multiplies it by
$q^{-2D(h)}$ with $D(h)=\tfrac12-h$. Because $2D(h)\in(0,1)$ on the admissible
range $0<h<\tfrac12$, the Laplacian acts on profile weights as a sum of
\emph{two} distinct homogeneous pieces, of degrees $b-1$ (radial) and
$b-2D(h)$ (axial). The radial degree is $b-1$, which is also the weight an
isotropic Laplacian carries in $q$; the axial degree equals it only at $h=0$,
and the two differ by $2h$.
The viscous operator therefore cannot be treated as a single homogeneous
degree-$(-2)$ object in the exponent ledger, and the geometric container
$|z|<q^{D}$ is what absorbs the mismatch.

Both mechanisms share one root: the kernel audits arithmetic that is reflected
into types and is silent about arithmetic that is not. We state precisely what
the artifact verifies (no declared axioms, no \ct{native\_decide}, no \ct{unsafe}, no
\ct{partial}, no \ct{implemented\_by}; exact exponent arithmetic; a
comparison theorem whose hypotheses are nearly empty) and what it does not (any
numerical size, any dimension, any physical realizability). The conclusion is a
verdict about statement design and verification scope, not about soundness: the
logical integrity of the artifact is robust, and the theorem it states is machine-checked against the statement it encodes. Kernel acceptance of the whole-space library was
reproduced on Linux (9,580 build jobs, exit~0, 817 modules; first in a virtual
machine, then from a pristine clone in continuous integration), after the pinned prerelease toolchain's \ct{lake} was found to
abort on current macOS; \ct{\#print axioms} on the four R3 theorems reports only
\ct{propext}, \ct{Classical.choice} and \ct{Quot.sound}, with no \ct{sorryAx}.
A build verifies terms rather than statements, so the statement was checked
separately: Comparator exports the challenge and the solution against the same
target list---both the whole-space and the periodic theorem among its
targets---replays the solution in a fresh kernel, and passes; an independent second kernel (\ct{nanoda}) accepts the same export as well (E75, \S7). No step of
the upstream proof is shown to be false. The stated whole-space result matches the
data and the conclusion of alternative (C) of the Clay statement---zero is an
admissible datum, and a smooth compactly supported force is stronger than
Clay's rapid-decay requirement---and it is silent on the unforced alternatives
(A) and (B). The periodic Navier--Stokes corollary is audited separately as alternative (D)
(audit log, E70), with its comparator bridge in both directions; its axiom
footprint, which E70 could only infer, was printed on a clean clone in continuous
integration and is the same three axioms, so (C) and (D) now rest on the same
kind of evidence (E73). The Euler
development in the same repository was outside this audit; its Comparator
challenge has since been executed and passes on both kernels, which checks the
two Euler theorems against their statements and reads no estimate inside their
proofs (E79). Those statements are the two declarations the repository's own
metadata designates as its Euler formalizations, and the challenge's solution
class is proved equivalent to the development's rather than assumed to match
(E80). E81 then compares them with the Euler paper's Theorem~1.1: the encoding is stronger than the paper in five respects and weaker in two---the challenge's $C^1$ limsup is the sum of the velocity and gradient suprema, while the paper asserts the gradient alone diverges, and the lifespan is taken over a Sobolev class not proved to contain the paper's smoothly evolving solution (E82); the development proves the gradient-only statement (\ct{maximalGradientNorm\_limsup}), its axiom footprint on a runner is the three permitted axioms with no \ct{sorryAx} (E84), and the Comparator target list does not contain it. Finally, the repository publishes the prose account it formalizes, and this audit now holds that account against the tree: the 166-page paper is mapped section by section onto the modules that realise it, ten sections and three appendices, each row labelled by what backs it---a citation in the source, an object match, or a family resemblance (E89). Five paper sections are named in the source at all, and eight of the thirteen rows rest on something weaker than a citation, so the map is section-level and not a clause-by-clause reconciliation.
\end{abstract}

\medskip
\noindent\textbf{Keywords:} formal verification, Navier--Stokes equations,
similarity ansatz, dependent type theory, proof auditing, statement design.

\noindent\textbf{Data availability:} commit \texttt{f9e8bc5} of
\texttt{openai/NavierStokesAndEuler}; audit log
\file{ns-audit/AUDIT\_LOG.md}, entries E0--E148. Every code citation below is a
\file{path:line} pair in that commit. Kernel acceptance was reproduced on the
\file{NavierStokes} library; see \S\ref{sec:limits}.

\tableofcontents

% =====================================================================
\section{Introduction}
\label{sec:intro}

\subsection{The artifact}

The repository under audit is a \Lean{}~4 development, built on Mathlib~4, of a
theorem about the three-dimensional incompressible Navier--Stokes equations. It
is large: $641{,}332$ lines in $2{,}659$ files, excluding the transitive
dependency closure (the frequently quoted figure of $28.6$ million lines is the
\ct{lean4export} closure including Mathlib, not this repository). It is also
unusually clean at the level of foundational hygiene: no \ct{axiom} declaration,
no \ct{native\_decide}, no \ct{unsafe}, no \ct{partial def}, no
\ct{@[implemented\_by]}, and no \ct{sorry} reachable from the proof roots. The
single \ct{sorry}-bearing module, \file{ComparatorChallenges/NavierStokes.lean},
is unreachable: a search for \ct{import ComparatorChallenges} over
\file{NavierStokes/}, \file{Euler/}, \file{NavierStokes.lean} and
\file{Euler.lean} returns nothing, and the module's own header states that
neither the proof root nor the submission imports it. The repository reports
acceptance by nine independent external checkers; we did not re-run them, apart from Comparator, which was run and passes on both of its challenges (E75, E79)
(\S\ref{sec:limits}).

This is, by any measure, a serious piece of proof engineering. The question this
paper asks is not whether it compiles. The question is what, precisely, the
compiled object certifies.

The audit is confined in depth to that whole-space forced result. The periodic
corollary is examined in \S\ref{sec:periodic} as alternative (D): its statement,
its derivation from the whole-space theorem, and its comparator bridge in both
directions. That examination is a reading of the statement and the architecture
of the proof, not an independent audit of the periodic development's own
machinery---the periodization lemmas, the periodic uniqueness theorem and the integration-by-parts layer beneath it are read rather than sampled (E83, E85, E86), though not re-proved---and the one empirical obligation it left, the axiom print, is
discharged in (E73). The Euler development in the same repository was outside
this paper until E79, which runs its Comparator challenge end to end: the
library builds from a clean clone (10,603 jobs, 50 min 11 s), both Euler
theorems are accepted by the builtin kernel and by nanoda, and their axiom
footprint is \ct{propext}, \ct{Classical.choice} and \ct{Quot.sound} with no
\ct{sorryAx}. That checks two proofs against two statements. E80 then checks the statements: the two are the declarations
\file{formalization.yaml:72--94} designates, both in \file{Euler/Solution.lean},
and the passage between the challenge's ordinary-function class and the
development's smooth--$L^2$ fields is a proved equivalence
(\ct{exists\_sobolevSolution\_iff}), with the decay and norm clauses bridged by
\ct{initialVelocityConditionDecay\_of\_compact} and the three norm identities.
No estimate inside the Euler development was read, so no mathematical claim
made here applies to it. E81 compares the two statements with the Euler paper's
own Theorem~1.1, clause by clause: the encoding is stronger than the paper in
five respects and weaker in two---the challenge's $C^1$ limsup adds the velocity
supremum to the gradient's, while the paper asserts only the gradient's
divergence, and the lifespan is defined over a class not proved to contain the
paper's smoothly evolving solution. The development proves the gradient-only
statement (\ct{maximalGradientNorm\_limsup}) and its axiom footprint is the three
permitted axioms (E84); the accepted statement does not ask for it.

\subsection{What the artifact proves}

The headline declaration is a five-line corollary of a longer construction.

\begin{lstlisting}[caption={(\file{NavierStokes/R3/ProblemStatement.lean:150}) The
statement the repository calls \ct{breakdownStatement}.},label={lst:breakdown}]
def breakdownStatement : Prop :=
  ∀ ν : ℝ, 0 < ν →
    ∃ u : VelocityField, ∃ p : PressureField, ∃ f : VelocityField, ∃ K : Set Space,
      CandidateProperties ν u p f K ∧ ¬ Nonempty (GlobalFiniteEnergySolution ν f)
\end{lstlisting}

The predicate \ct{CandidateProperties} (\file{NavierStokes/R3/ProblemStatement.lean:92})
bundles twelve fields; the ones that matter here are
\ct{velocity\_smooth} and \ct{pressure\_smooth} ($C^\infty$ on the pre-singular
domain), \ct{speed\_unbounded} ($\limsup_{t\uparrow 1}\norm{u(t)}_\infty=\infty$),
\ct{zero\_initial\_velocity} ($u(0,\cdot)\equiv0$), \ct{divergence\_free},
\ct{force\_smooth} and \ct{force\_support} (smooth, with time support in
$t>0$), compact spatial support for $u$ and $p$, and the equation itself:
\[
  \ct{navier\_stokes}:\qquad
  \partial_t u+(u\cdot\nabla)u-\nu\Delta u+\nabla p=f
  \quad\text{on }(0,1)\times\R^3 .
\]
The refined form actually proved adds a rest interval:
\ct{theorem\_1\_1\_with\_initial\_rest} (\file{NavierStokes/R3/Theorem.lean:26})
produces $u\equiv p\equiv0$ for $|t|\le 3/8$, and
\ct{theorem\_1\_1} (\file{NavierStokes/R3/Theorem.lean:46}) is the
constant-free restatement. In prose:

\begin{quote}
For every $\nu>0$ there exist $u,p,f$ and a compact $K\subset\R^3$ such that
$u,p$ are smooth and divergence-free on the pre-singular domain with spatial
support in $K$, $f$ is smooth with time support in $t>0$, the forced
Navier--Stokes system holds identically, $u\equiv p\equiv0$ on $|t|\le3/8$ (so
$u(0,\cdot)=0$), $\sup_{t\in[0,1)}\tfrac12\norm{u(t)}_2^2<\infty$,
$\limsup_{t\uparrow1}\norm{u(t)}_\infty=+\infty$, and \emph{no} global smooth
finite-energy solution of the same forced problem starting from the zero datum
exists.
\end{quote}

\medskip
\noindent
That proposition is the repository's Theorem~1.1, not Clay's alternative~(C) as
written. The bridge is the Comparator challenge, and E76 checks it clause by
clause against Fefferman's text: each numbered condition of (C) and (D) --- the
equation, incompressibility, the initial condition, the decay bounds (4), (5),
(8), (9), the smoothness and energy clauses (6), (7), (10), (11), and the
pressure-periodicity erratum --- maps onto a named field of
\ct{InitialVelocityConditionDecay}, \ct{ForceConditionDecay},
\ct{NavierStokesExistenceAndSmoothnessRn} and their periodic counterparts, with
no clause missing, weakened, or strengthened. Two reading conventions are named
rather than hidden: $C^\infty$ on the closed half-space is \ct{ContDiffOn} on
$\R^3\times[0,\infty)$, and the multi-index bounds are carried as iterated
Fr\'echet operator norms, which dominate every $|\partial^\alpha|$ up to a
constant that the existential absorbs.

\subsection{Two halves, of very different cost}
\label{sec:twohalves}

The statement is a conjunction, and the two conjuncts have almost nothing to do
with each other. Separating them is the first step of any honest accounting.

\begin{proposition}[Naming the force is free; smoothing it is not]
\label{prop:free}
Let $u\in C^\infty([0,1)\times\R^3;\R^3)$ be divergence-free on
$[0,1)\times\R^3$, compactly supported in $x$ uniformly in $t$, with
$u\equiv0$ for $t\le3/8$, with
$\limsup_{t\uparrow1}\norm{u(t)}_\infty=\infty$, and with
$\sup_{t\in[0,1)}\norm{u(t)}_2<\infty$. Put $p:=0$, let $K$ be the closure of
$\bigcup_{t\in[0,1)}\supp\bigl(u(t,\cdot)\bigr)$, and set
\[
  f:=\partial_t u+(u\cdot\nabla)u-\nu\Delta u
  \qquad\text{on }[0,1)\times\R^3 .
\]
Then
\begin{enumerate}
\item[(i)] $f$ is $C^\infty$ on $[0,1)\times\R^3$, satisfies
$\supp f(t,\cdot)\subseteq K$ for every $t\in[0,1)$, and is identically zero
for $t\le3/8$;
\item[(ii)] $\ct{navier\_stokes}$ holds identically on $(0,1)\times\R^3$;
\item[(iii)] every field of \ct{CandidateProperties} except
\ct{force\_smooth} and \ct{force\_support} holds of the slab data.
\end{enumerate}
What does \emph{not} follow is the pair of global fields: a $C^\infty$
extension of $f$ to $\R^4$, and the compact-support condition stated for a
function on $\R^4$.
\end{proposition}

\begin{proof}
For (i): $f$ is a finite combination of $C^\infty$ operations applied to $u$
on an open slab on which every term is smooth, so it is $C^\infty$ there; its
spatial support is inherited from $u$, and it vanishes near $t=0$ because $u$
does. For (ii): the equation is an identity by the definition of $f$. For
(iii): substitute (i) and (ii) into the field list. No step extends $f$ to
$t=1$: $\ct{force\_smooth}$ is a statement about $\R^4$, and $f$ has been
defined on a proper subset of it.
\end{proof}

Proposition~\ref{prop:free} isolates what is genuinely free: no fluid equation
is solved, a field is chosen, and its residual is named ``force''. The
repository does not conceal this: \ct{originalResidual A v p := navierStokesResidual
(velocity A v) p} defines the force as the residual of the constructed fields,
and \ct{force\_nonzero\_of\_no\_global\_solution}
(\file{NavierStokes/R3/CandidateBreakdown.lean:53}) records, as a lemma, that
any force in this class which excludes a global solution must be nonzero
somewhere at positive time:
\[
  \ct{CompactPositiveTimeSupport}(f)\ \wedge\
  \neg\,\exists\,\text{global finite-energy solution}(\nu,f)
  \ \Longrightarrow\
  \exists t>0\ \exists x,\ f(t,x)\neq0 .
\]
That is not an admission of a gap; it is a formal lower bound on the force,
and it is the sharpest available statement that the result is \emph{forced}
rather than spontaneous.

The distance between Proposition~\ref{prop:free} and the theorem is not a gap in
the argument; it is a second obligation, and it is the expensive one. Because
\ct{speed\_unbounded} forces $\partial_t u$ to be singular at $t=1$, the residual
is a sum of singular terms whose combination has to cancel to all orders there.
Producing a field that blows up is easy; producing one whose residual does not
is the construction. The repository discharges that obligation in the force
assembly rather than through the identity: \ct{of\_localized\_fields}
(\file{NavierStokes/R3/ActualCandidate.lean:76}) takes
$h_f:\ct{ContDiff}\ \R\ \infty\ f$ as a \emph{hypothesis} and passes it
straight into the \ct{force\_smooth} field. The ledger therefore has three
entries, not two.

\begin{remark}[What is \emph{not} free]
\label{rem:notfree}
Two conjuncts carry the cost. The first is \ct{force\_smooth}: a residual
defined on the pre-singular slab must be extended smoothly across the time at
which its own $\partial_t u$ term is singular. The second is
$\neg\,\mathsf{Nonempty}(\ct{GlobalFiniteEnergySolution}\ \nu\ f)$, a genuine
uniqueness theorem. It is proved by
\ct{classical\_uniqueness\_on\_Icc} (\file{NavierStokes/R3/WholeSpaceUniqueness.lean:30}),
whose competitor class imposes smoothness, divergence-freeness, agreement of
residuals, and uniform finite energy---and nothing else: no pressure norm, no
spatial decay of $p-q$ at infinity, no support, growth, or derivative
hypothesis on the competitor. The pressure is eliminated through a Riesz
pairing, and the Gronwall constant is uniform in the localization radius by
binder order ($\exists D\ge0,\ \forall R\ge1$). A tautological force does not
buy any of that; a wrong force class would destroy it. So the ledger reads:
\emph{naming} $=$ free; \emph{smoothing} $=$ the construction;
\emph{comparison} $=$ the mathematics.
\end{remark}

\subsection{Two structural findings}

Section~\ref{sec:constants} isolates the \emph{Non-Evaluated Constant Barrier}:
the analytic chain is quantitatively vacuous at the level of constants, and
this is invisible to the kernel for a type-theoretic reason, not an accidental
one. Section~\ref{sec:anisotropy} isolates \emph{Anisotropic Viscous Scaling}:
the repository's own pullback identities show that the Laplacian is a sum of
two homogeneous pieces in the similarity weight---the radial group of degree
$1$, the axial group of degree $2D(h)=1-2h\in(0,1)$---so that no
admissible member of the family is isotropically self-similar.

The two findings are instances of one principle, stated in
\S\ref{sec:index-vs-term}: \emph{the kernel audits arithmetic that is reflected
into types, and is silent about arithmetic that is not.} Constants in
existential binders and real exponents inside definitions are not reflected.

\subsection{Relation to the macroscopic caveats}

Public discussion of this repository has concentrated on two features: the zero
initial datum, and the forcing term. Both are real, both are statement-level,
and neither is the subject of this paper. We take them as given and ask a
different question: given that statement, what does the proof \emph{certify}
about it? The answer is not ``everything the compiler checked''. The
distinction matters because ``machine-checked'' is routinely read as
``numerically and dimensionally reliable'', and \S\ref{sec:constants} shows that
reading is false here for a structural reason that will recur in any large
formalization of a quantitative PDE estimate.

\subsection{Method and reproducibility}
\label{sec:method}

The audit proceeds by reading proof terms and binder sequences, not by
re-elaborating them. For each claim we record the declaring file and line, the
binder order where it matters, the informal statement in standard notation, and
a dimensional check. Where a matrix of review tasks named a file or declaration
that does not exist at \texttt{f9e8bc5}, the log records the missing premise and
points at the module that actually carries the content, rather than inventing a
lemma. Of the thirty-five tasks in the final review matrix, nineteen were
answered from declarations that exist under the names given; sixteen required
premise correction (audit log, E46 and E47).

Counts are re-derived from the source text rather than from earlier totals: the
denominator of the result ledger was corrected from 74 to 73 when the per-kind
breakdown behind it was found to have counted 38 lemmas where the text has 37
(E135).

\subsection{Three kinds of claim, and which ones this audit makes}
\label{sec:categories}

The results here are easy to over-read, so we separate three kinds of assertion
and say which kind each one is.

\begin{enumerate}
\item \textbf{A property of the formalized code.} Something true of the source at
\texttt{f9e8bc5}: a binder appears in this order, an identity is closed by
\texttt{rfl}, a record field is inhabited by \texttt{Classical.choose}, a module
is not imported by the proof root. Claims of this kind are settled by reading,
and every one in this paper carries a \file{path:line} pair.

\item \textbf{A limitation on what the formalization establishes.} A claim about
the \emph{type} rather than the \emph{term}: no numeral is recovered from an
existential constant, no unit is weighed, no homogeneity class is asserted. Such
a limit can hold of a development whose theorem is also true, and it is the kind
of claim both findings make.

\item \textbf{A mathematical error that invalidates a proof step.} An estimate
that is false, an exponent that is wrong, an inference that does not follow. This
is the only kind that bears on whether the theorem is true, and it requires
naming the step, its \file{path:line}, and the inequality that fails.
\end{enumerate}

Sections~\ref{sec:constants} and the anisotropy section establish claims of the
first kind and draw consequences of the second; the three interface audits record
identities of the first kind. No claim of the third kind is made anywhere in this
paper. The one place the audit briefly made such a claim --- that the anisotropic
weights admitted no order-by-order expansion --- is recorded and withdrawn in the
log (E63), because the source contradicted it. That withdrawal is the standard
the remaining claims are held to: a category-three assertion survives only if the
defective step can be exhibited.

A fourth outcome has since been forced by the log, and it is not one of the three.
When a paper clause quantifies over the \emph{derivation} rather than over the
fields---how many derivatives of each source an estimate consumes---the object
language has no predicate for it. Lean terms are finite, so ``finitely many input
derivatives'' cannot fail, and there is no proposition to state. The third clause
of Lemma~9.7 is of this kind (E133). The audit records such a clause as
\emph{unstatable} rather than missing, and does not count it as evidence of a gap;
it is distinct from a query that used the wrong name (E94) and from a component
stated in place of a statement (E129).

Two rows run the other way and belong in the same paragraph, because they are the
same observation with the sign reversed: the development is \emph{stronger} than
the prose where it states the Fr\'echet derivative of the torus inverse rather than
the equation for its series (E132), and where it drops the flat error term of the
residual estimate instead of bounding it (E134). Neither is evidence against the
paper. They are recorded so that a later reader does not mistake the direction of
the divergence.

% =====================================================================
\section{Preliminaries: notation and the two spaces}
\label{sec:prelim}

We write $u$ for velocity, $p$ for pressure, $f$ for the forcing, $\nu>0$ for
kinematic viscosity, $\omega=\nabla\times u$ for vorticity, and
$w:=u-v$ for the difference of two velocity fields. $\tau:=1-t$ is the time to
the singular time, $r=|x'|$ the cylindrical radius, $s:=r^2/2$, and $z:=x_3$
the axial coordinate. $\T^3$ denotes the unit torus and
$\ct{innerCube}(r):=\{x\in\R^3:\abs{x_i}<1-r\ \forall i\}$
(\file{NavierStokes/PeriodicLocalization.lean:213}); at $r=1/4$ this is
$(-\tfrac34,\tfrac34)^3$.

Two facts about the whole-space/torus interface are worth recording because they
were the target of an earlier hypothesis that the audit refuted. First, the
whole-space development is primary and the torus is periodized \emph{from} it,
by a lattice sum of the cut whole-space pressure, of the form
$p_{\T^3}(t,x)=\sum_{k\in\mathbb{Z}^3}\chi(x+k)\,p(t,x+k)$, whose local finiteness
is proved rather than assumed. Second, agreement between the periodic and cut
residuals is an exact jet identity on the inner cube,
\[
  \ct{periodicResidual\_jets\_eq\_cut}:\quad
  \partial^{\alpha}\ct{periodicResidual}\,A\,p\,(t,x)
  =\partial^{\alpha}\ct{cutResidual}\,A\,p\,(t,x)
  \qquad \bigl(x\in\ct{innerCube}(\tfrac14)\bigr),
\]
(\file{NavierStokes/PeriodicResidualLimits.lean:138}), obtained from a germ
equality of the two residual functions rather than from a decay hypothesis. The
support of the construction has radius at most $1/4$ while the inner cube has
margin $1/2$, so the lattice sum has no overlap and no spectral truncation is
involved.

The relevance for this paper is narrow but real: the pressure entering the
comparison is determined only up to a constant (the Riesz symbol is applied to
the data and the constant is annihilated by $\nabla\cdot w=0$), so no pressure
norm is ever assumed. That is a strength of the artifact, and it is why the
comparison survives the statement-design objections. It is also why the
\emph{quantitative} content has to be examined separately: with the pressure
eliminated, everything that remains is a chain of inequalities, and that chain
is where \S\ref{sec:constants} bites.

\subsection{Alternative (D) and the pressure-periodicity clause}
\label{sec:pressure-periodic}

The repository also formalizes a periodic corollary, corresponding to
alternative (D) rather than (C). Its candidate \emph{and competitor} structures
carry a clause with no whole-space counterpart
(\file{NavierStokes/PeriodicPaperTheorem.lean:36}, \file{:54}):
\begin{lstlisting}
pressure_periodic : UnitSpatialPeriodsOn (Ico 0 1) p
\end{lstlisting}
Read against a statement that imposes periodicity on the velocity alone, this
would narrow the competitor class, and narrowing that class is exactly how a
breakdown claim can be made to look easier than it is. That is the objection,
and it does not apply here: the official problem description carries an erratum
making explicit that in the periodic case equation~(8) requires both the
velocity and the pressure to be periodic. The repository's clause is the
published hypothesis, not an addition to it.

The clause is inert for everything in this paper, which concerns alternative
(C) on $\R^3$ and never invokes pressure periodicity; the whole-space
comparison eliminates the pressure through a Riesz pairing instead
(\S\ref{sec:interfaces}). We record one limitation in the same breath as the
claim: the erratum was not read at first hand for this audit. The attribution
rests on the repository's own reference and on secondary descriptions, so the
clause is correctly attributed rather than independently verified.

\subsection{The periodic corollary, audited as alternative (D)}
\label{sec:periodic}

The repository also proves a periodic statement, and it is worth separating from
the whole-space result because its quantifiers differ in a direction that favours
the theorem. The corollary asserts: for every $\nu>0$ there exist a velocity, a
pressure, a force and a compact set $K$ in the interior of the fundamental cube,
such that the fields are smooth and periodic before time one, the initial
velocity vanishes, the residual equals the force, the speed is unbounded at one,
and no global smooth periodic solution exists for that same force.

Two features of the statement deserve emphasis, because each is stronger than its
whole-space counterpart rather than weaker.

First, the excluded competitor class imposes no energy bound. A competing
solution must be smooth, periodic in space (velocity and pressure),
divergence-free, vanish at time zero, and satisfy the equation for all $t>0$. It
need not have bounded kinetic energy, a normalised pressure, or any decay or
growth condition. Excluding a larger class is the stronger theorem, so the
periodic result is not bought by weakening the competitor; if anything the
whole-space alternative (C) is the more cautious encoding of the two.

Second, the passage from the whole-space theorem is a genuine similarity, not a
relabelling. The clock $l^2(t-1)+1$ fixes the singular time at exactly one, space
is compressed by $1/l$, the force scales by $l^3$ and the pressure by $l^2$, and
the viscosity is untouched. That the compressed object still solves the equation
with the transformed force is checked at the level of the residual, and the
pre-initial history introduced by the clock shift is removed before the lattice
sum. The choice of $l$ is what places the supports inside the interior of the
fundamental cube, after which periodization is applied to the fields with their
post-one extensions discarded.

The adapter to the comparator runs in both directions: a candidate yields the
comparator's exact existential quantifiers with zero datum and the argument order
swapped, and a comparator solution converts back into the excluded class. Neither
direction rescales viscosity or time.

One statement-design remark, disclosed by the source itself: ``support'' here is
not compact support on $\R^3$, which a nonzero periodic lift cannot have. The
record intersects the support of each time slice with the fundamental cube and
places that intersection inside $K$. Read as a claim about all of $\R^3$ the
clause would be false; read against the torus, it is the intended one.

The limitation recorded in E70 was empirical, and it is now discharged: on a clean
clone, the axiom probe printed the periodic theorem, its paper corollary and both
comparator adapters, and each depends only on \ct{propext}, \ct{Classical.choice}
and \ct{Quot.sound} (E73). The four whole-space statements had already been
printed that way; these four joined them, and the same run printed the
whole-space comparator theorem as well.

% =====================================================================
\section{The Non-Evaluated Constant Barrier}
\label{sec:constants}

\subsection{Three exhibits}

\begin{example}[A Sobolev constant inherited from Mathlib]
\label{ex:sobolev}
\file{NavierStokes/R3/WeightedSobolev.lean:25} defines
\begin{lstlisting}
def sobolevConstant : ℝ :=
  (eLpNormLESNormFDerivOfEqInnerConst (volume : Measure Space) 2 : ℝ)
\end{lstlisting}
and \file{:42} uses it in the sharp Gagliardo--Nirenberg--Sobolev inequality for
$C^1_c$ fields,
\[
  \norm{f}_{L^6(\R^3)}\ \le\ \ct{sobolevConstant}\cdot
  \norm{\nabla f}_{L^2(\R^3)} .
\]
The constant is \ct{noncomputable}: it is whatever Mathlib's proof of the
inequality produced. A derived constant, \ct{weightedSobolevConstant :=
4 * sobolevConstant + 1} (\file{:210}), is likewise opaque.
\end{example}

\begin{example}[A witness extracted from an existence proof]
\label{ex:choose}
\file{NavierStokes/R3/LocalizedFluxEstimates.lean:176} proves
\[
  \exists C>0,\ \forall x\in\R^3,\quad
  \norm{D^2\ct{baseWeight}(x)}\le C,
\]
by compactness of the support and continuity of the second iterated derivative,
taking $C=\max(1,C_0)$. Line \file{:184} then \emph{names} the witness:
\begin{lstlisting}
def weightSecondDerivativeConstant : ℝ :=
  Classical.choose exists_weight_second_derivative_bound
\end{lstlisting}
with \file{:186} and \file{:190} recording only positivity and the bound. The
term is definitionally opaque: it is not a reducible definition, it has no
computable content, and it is not characterized by any equation beyond the
inequality it was chosen to satisfy.
\end{example}

\begin{example}[A pressure functional bounded by an existential constant]
\label{ex:pressure}
\file{NavierStokes/R3/PressureFunctionals.lean:169} establishes
\[
  \exists C\ge0,\ \forall \psi\in\mathcal{S}(\R^3;\mathbb{C}),\quad
  \bigl\lvert \ct{averagedPressureDifferenceValue}\,W_0\,W_1\,G\,k\,\psi
  \bigr\rvert
  \ \le\ C\,\sqrt{\ct{fourierHNormSq}\ 3\ \psi},
\]
and \file{:279} restates it for the linear functional actually consumed by the
flux estimate. The witness is assembled from an explicit expression in the data,
\[
  M=3\sqrt{\ct{l2Sq}\,W_0}+\sqrt{\ct{l2Sq}\,W_1}
   +\sum_i\int\norm{G_{ki}}
   +\sum_{i,j}\int\norm{G_{ij}},
\]
\emph{times} an operator-norm constant $B$ supplied by
\ct{exists\_uniform\_test\_bound}. So the dependence on the data is explicit;
the numerical size is not. This is the most favourable of the three cases, and
it still does not evaluate.
\end{example}

\subsection{Why the kernel cannot see it: indices versus terms}
\label{sec:index-vs-term}

The three examples share a mechanism, and it is worth stating in
type-theoretic rather than rhetorical terms.

\paragraph{Indices are checked.} In dependent type theory a family such as
$\ct{Lp}\ \mathbb{C}\ p$ with $p:\mathbb{R}_{\ge0}\cup\{\infty\}$, or
$\ct{ContDiff}\ \R\ n$ with $n:\mathbb{N}\cup\{\infty\}$, is a \emph{type
family}. A term of type $\ct{Lp}\ \mathbb{C}\ 6$ is not a term of type
$\ct{Lp}\ \mathbb{C}\ 4$. To use one where the other is expected, the
elaborator must solve the unification problem $6\stackrel{?}{=}4$, which fails.
Hence: \emph{any exponent reflected into an index is checked by unification, and
changing it by a nonzero amount is a type error.} For such exponents the kernel
is a complete decision procedure---\emph{relative to the declared type family}.
It decides whether a term inhabits the type it was given, not whether that type
is the one the physics requires; a development that consistently encodes the
wrong exponent in an index elaborates exactly as smoothly as a correct one. The
guarantee is internal consistency with the specification, not correspondence
between the specification and the problem.

\paragraph{Terms are only inhabited.} By contrast
\[
  \exists C:\R,\ 0\le C\ \wedge\ \forall x,\ \norm{Tx}\le C\norm{Sx}
  \tag{$\star$}
\]
is a proposition. A proof term is checked against $(\star)$ by verifying that
the witness inhabits it. There is no second, independent specification of $C$
against which the witness could be compared; the proposition \emph{is} the
specification. And once $C$ is produced by \ct{Classical.choose} from an
existence proof, it is not merely unknown but \emph{inaccessible}: not
reducible, not computable (\ct{noncomputable}), not inspectable by
\ct{norm\_num} or \ct{decide}, and with no defining equation.

\paragraph{Consequence.} An error in an exponent that reaches an index becomes a
compile failure. An error in a multiplicative constant does not: the
proposition $(\star)$ remains true for $C$ replaced by $10^{6}C$, and every use
in the repository is monotone in $C$ (each occurrence is of the form
``$\le C\cdot(\text{norm})$'' with $C$ in positive position), so the inflated
witness inhabits the same statements. The artifact therefore cannot distinguish
a correct constant from an inflated one. We emphasize that this monotonicity
observation is a meta-argument about the proof scripts, not a machine-checked
statement; it is the reason the audit had to verify dimensions by hand
(audit log, E10.7).

\subsection{Epistemic, not logical}
\label{sec:epistemic}

It is essential not to overstate the finding. The barrier does not create
unsoundness and does not threaten the truth of the theorem.

\begin{proposition}[The barrier cannot falsify the headline]
\label{prop:constantfree}
The type of \ct{theorem\_1\_1} contains no numerical constant: it is
$\forall\nu>0,\ \exists u\,\exists p\,\exists f\,\exists K,\
\ct{CandidateProperties}\ \nu\,u\,p\,f\,K\wedge\neg\mathsf{Nonempty}(\dots)$.
Consequently no assignment of values to the constants appearing in auxiliary
lemmas can change the truth value of the statement proved.
\end{proposition}

\begin{proof}
Immediate from the shape of the type in Listing~\ref{lst:breakdown} and
\file{NavierStokes/R3/Theorem.lean:46}. No constant occurs in the binders or in
the body except through \ct{CandidateProperties}, whose fields are qualitative
(smoothness, support, divergence-freeness, an identity, an unboundedness
statement, and a negated existential).
\end{proof}

What the barrier does foreclose is a specific and frequently assumed class of
claims:

\begin{center}
\begin{tabular}{@{}p{0.52\textwidth}p{0.40\textwidth}@{}}
\toprule
Claim & Status under the artifact \\
\midrule
The proof yields an explicit bound & \textbf{No}: every constant is
$\exists$-bound or \ct{choose}n. \\
The bound is dimensionally consistent & \textbf{Not machine-checked}:
dimensions are a semantic layer absent from the type theory. \\
The constant is small enough for the next step & \textbf{No}: no step
instantiates a constant by a numeral. \\
The exponent bookkeeping is correct & \textbf{Yes}, where the exponent reaches
a type index or a reducible definition (\S\ref{sec:ledger}). \\
The theorem itself holds & \textbf{Yes}, subject to kernel acceptance. \\
\bottomrule
\end{tabular}
\end{center}

\subsection{The contrast case: exponents that \emph{are} checked}
\label{sec:ledger}

It would be wrong to conclude that the formalization checks nothing
quantitative. The exponent ledger does real work, and it does it precisely
because its objects are reducible definitions rather than existential witnesses.

\file{NavierStokes/ExponentLedger.lean} fixes the manuscript's residual-order
budget. Its header is explicit that it ``checks the arithmetic of those
assignments, conditional on the estimates being valid'' and that it ``does not
define the analytic classes, construct a correction, or prove an estimate for a
PDE''. Within that scope the results are exact. The particular-correction gain is
defined by a nested minimum,
\[
  \ct{particularGain}(\sigma,\kappa)
  =\min\bigl(\min(\min(\tfrac12-3\kappa,\ \tfrac12-\kappa),\ \tfrac12+\sigma-\kappa),\ \tfrac25\bigr),
\]
and \ct{particular\_gain\_eq} proves
\[
  \ct{particularGain}(\sigma,\kappa)=\tfrac25
  \qquad\text{for }\ \sigma\ge\tfrac15,\ \kappa\le10^{-5},
\]
by \ct{linarith} on exact rationals. The signed correction's gain is
$\ct{signedGain}(\sigma,\kappa)=\tfrac25-\kappa$ under the same hypotheses, so
the binding gain is $\tfrac25-\kappa$ and the first listed entry
$\tfrac12-4\kappa$ is slack by $\tfrac1{10}-3\kappa$.

This is the sharpest possible contrast with \S\ref{sec:index-vs-term}. The gain
is a \emph{definition} of a real number built from rationals; equality with a
numeral is decidable by linear arithmetic, and the kernel checks it. The Sobolev
constant is a \emph{witness} of an existential; no numeral is available, and the
kernel checks only that the witness exists.

Two further consequences of the ledger deserve record. First, $\kappa$ is
\emph{never defined} in the repository: every theorem is stated for
$\kappa\le1/100000$, and \ct{manuscript\_kappa\_admissible} merely checks the
numeral $10^{-5}$ against that bound. The artifact therefore verifies a
\emph{range}, not the manuscript's choice; \ct{smaller\_kappa\_admissible}
records that $10^{-6}$ works as well. Second, the most restrictive numerical
margin is the $\tfrac{17}{100}$ bar-residual bound, whose exact threshold is
$\kappa<\tfrac1{200}$ (\ct{old\_difference\_bar\_margin\_iff}), three orders of
magnitude above the value used. The ledger is thus not tight, and its
correctness is a statement about consistency, not optimality.

\subsection{Consequence for the ``fully verified'' reading}

The natural reading of ``machine-checked'' is that the quantitative content of
the paper has been checked. On this artifact that reading fails, and it fails
structurally: not because a constant is wrong, but because the proof assistant
has no mechanism by which a constant could be \emph{right} in a checkable sense.
The repository's formalization is, at the quantitative level, a certificate of
\emph{consistency of the exponent algebra and the implication structure}. That is
a real and non-trivial achievement---it is what catches an off-by-one in a
derivative count, and such errors are the common failure mode in this
literature---but it is not a certificate of numerical or dimensional
correctness.

% =====================================================================
\section{Anisotropic Viscous Scaling under the Generalized Cone}
\label{sec:anisotropy}

\subsection{The coordinates}

The similarity coordinates are those of the manuscript, formalized in
\file{NavierStokes/SimilarityCoordinates.lean}. With
$a=2h\in(0,1)$ and $\tau=1-t>0$, the quantity $q=q(\tau,z)$ is the unique
positive solution of
\begin{equation}
  q-z^{2}q^{2h}=\tau ,
  \label{eq:q}
\end{equation}
(\ct{forwardScalar}, \ct{coordinateQ\_spec}); the inner coordinates are
\begin{equation}
  \eta:=\frac{z}{q^{D(h)}},\qquad X:=\frac{s}{q},
  \qquad D(h):=\frac12-h,
  \label{eq:inner}
\end{equation}
matching \ct{coordinateEta} (\file{SimilarityCoordinates.lean:340}) and
\ct{coordinateX} (\file{:500}), with $D$ at
\file{CoordinateAlgebra.lean:19}. Eliminating $\eta$ from \eqref{eq:q} using
$z=q^{D}\eta$ gives the identity
$q(1-\eta^{2})=q-z^{2}q^{2h}=\tau$, i.e.\ $\tau=q\,d(\eta)$ with
$d(\eta)=1-\eta^{2}$; and since $a=2h$,
\[
  (q^{D})^{2}q^{2h}=q^{2D+2h}=q ,
\]
so that \eqref{eq:q} is exactly
\begin{equation}
  q\,d(\eta)=\tau .
  \label{eq:qeta}
\end{equation}

The geometric content is in \eqref{eq:inner}: the level set $\eta=\pm1$, i.e.\
\begin{equation}
  \abs{z}=q^{D(h)},
  \label{eq:cone}
\end{equation}
is the boundary of the \emph{generalized cone}. For $h\in(0,\tfrac12)$ one has
$D(h)\in(0,\tfrac12)$, so \eqref{eq:cone} is a paraboloid-like surface
$|z|\sim\tau^{D}$ with $D=h$-dependent opening; the classical parabolic case is
the opposite endpoint $h=0$, where $D(0)=\tfrac12$ and $\eta=z/\sqrt{q}$
recovers the ordinary parabolic variable. The admissible range $h\in(0,\tfrac12)$
excludes that endpoint, and at $h=\tfrac12$ one has $D=0$, so the axial scale
ceases to contract. The scaling
group is
\begin{equation}
  (\tau,z,s,q)\ \longmapsto\
  (Q\tau,\ Q^{D(h)}z,\ Q\,s,\ Q\,q),
  \label{eq:scaling}
\end{equation}
which leaves $\eta$ and $X$ invariant. That this group acts by \emph{proved}
homogeneity rather than by postulation is the content of
\ct{forwardScalar\_scale}, \ct{coordinateQ\_scale\_h},
\ct{coordinateEta\_scale\_h} and \ct{coordinateX\_scale\_h}
(\file{NavierStokes/SimilarityHomogeneity.lean}); the last three derive the
scaling of the coordinate from the defining equation \eqref{eq:q} and
uniqueness of its positive branch, not from an assumed ansatz.

\subsection{The pullback and its four exact derivative identities}

\begin{definition}[Profile pullback]
\label{def:pullback}
For $h\in(0,\tfrac12)$, $b\in\R$ and an inner profile
$f:\R^2\to\R$, the physical profile is
\[
  \bigl(\Ph_{h,b}f\bigr)(t,s,z):=q^{b}\,f\bigl(X,\eta\bigr),
  \qquad (X,\eta)\ \text{as in }\eqref{eq:inner},
\]
(\file{NavierStokes/SimilarityProfile.lean:45}). The physical partial
derivatives are genuine Fréchet derivatives along the coordinate directions:
$\partial_tF:=DF\,(1,0,0)$, $\partial_sF:=DF\,(0,1,0)$,
$\partial_zF:=DF\,(0,0,1)$ (\file{:48}--\file{:50}).
\end{definition}

\begin{lemma}[Exact weight transport]
\label{lem:pullback}
Let $h\in(0,\tfrac12)$, $b\in\R$, $f\in C^{2}$ near $(X,\eta)$, and
$p=(t,s,z)$ with $t<1$. Then
\begin{align}
  \partial_s\bigl(\Ph_{h,b}f\bigr) &= \Ph_{h,b-1}\bigl(\partial_X f\bigr),
  \label{eq:ds}\\[2pt]
  \partial_z\bigl(\Ph_{h,b}f\bigr) &= \Ph_{h,b-D(h)}\bigl(Z_{h,b}f\bigr),
  \label{eq:dz}\\[2pt]
  \partial_s^{2}\bigl(\Ph_{h,b}f\bigr) &= \Ph_{h,b-2}\bigl(\partial_X^{2}f\bigr),
  \label{eq:dss}\\[2pt]
  \partial_z^{2}\bigl(\Ph_{h,b}f\bigr) &=
  \Ph_{h,b-2D(h)}\bigl(Z_{h,b-D(h)}Z_{h,b}f\bigr),
  \label{eq:dzz}
\end{align}
where $Z_{h,b}$ is the explicit first-order coefficient
\[
  Z_{h,b}f=\frac{2\eta\,b\,f+d(\eta)\,\partial_\eta f-2\eta X\,\partial_X f}
              {L(h,\eta)},
  \qquad L(h,\eta)=1-2h\eta^{2},
\]
(\file{CoordinateAlgebra.lean:21} and \file{:184}).
\end{lemma}

\begin{proof}[Proof]
Equations \eqref{eq:ds} and \eqref{eq:dz} are
\ct{partialS\_pullback} (\file{SimilarityProfile.lean:224}) and
\ct{partialZ\_pullback} (\file{:233}); \eqref{eq:dss} and \eqref{eq:dzz} are
\ct{partialS\_partialS\_pullback} (\file{:279}) and
\ct{partialZ\_partialZ\_pullback} (\file{:296}). Each is proved by
differentiating the pullback along the relevant coordinate using the chain and
product rules, then identifying the derivative with the stated pullback of the
corresponding inner expression; the chain coefficients are the explicit
rational functions \ct{qTime}, \ct{etaTime}, \ct{xTime}, \ct{qAxial},
\ct{etaAxial}, \ct{xAxial} of \file{CoordinateAlgebra.lean:138}--\file{:143},
and the identifications are \ct{time\_chain\_coefficient},
\ct{axial\_chain\_coefficient}. The second-order cases require the local
$C^2$ hypothesis to differentiate the first-order identity in a neighbourhood,
which is why they are stated with \ct{ContDiffAt} rather than pointwise
differentiability.
\end{proof}

\begin{remark}
The asymmetry between \eqref{eq:dss} and \eqref{eq:dzz} is the whole point:
$\partial_s$ transports the weight by exactly $-1$, while $\partial_z$
transports it by $-D(h)$. The former is a consequence of $X=s/q$; the latter of
$\eta=z/q^{D}$, and it is the cone's opening exponent that appears.
\end{remark}

\subsection{The Laplacian is a sum of two homogeneous weights}

In the axisymmetric variables $(s,z)$ the Laplacian is
\begin{equation}
  \Delta_{\mathrm{ax}}
  = 2s\,\partial_s^{2}+2\,\partial_s+\partial_z^{2},
  \label{eq:lap}
\end{equation}
(the factor $2$ is the Jacobian of $s=r^{2}/2$; only the \emph{weights} matter
below). Applying Lemma~\ref{lem:pullback} termwise requires one care: the
multiplicative factor $s=qX$ in front of $\partial_s^{2}$ is not constant in
the inner variables, so it travels with the derivative and consumes one power
of $q$.

\begin{proposition}[Two weights, not one]
\label{prop:three}
For $h\in(0,\tfrac12)$, $b\in\R$ and $f\in C^{2}$,
\begin{equation}
  \Delta_{\mathrm{ax}}\bigl(\Ph_{h,b}f\bigr)
  =
  \Ph_{h,b-1}\bigl(2X\,\partial_X^{2}f+2\,\partial_X f\bigr)
  +\Ph_{h,b-2D(h)}\bigl(Z_{h,b-D(h)}Z_{h,b}f\bigr).
  \label{eq:laplace-weights}
\end{equation}
\end{proposition}

\begin{proof}
By \eqref{eq:dss}, $\partial_s^{2}(\Ph_{h,b}f)=\Ph_{h,b-2}(\partial_X^{2}f)
=q^{b-2}\,\partial_X^{2}f$; multiplying by $2s=2qX$ gives
$2q^{b-1}X\,\partial_X^{2}f=\Ph_{h,b-1}(2X\,\partial_X^{2}f)$. By
\eqref{eq:ds}, $2\partial_s(\Ph_{h,b}f)=\Ph_{h,b-1}(2\,\partial_X f)$. By
\eqref{eq:dzz}, the axial term is $\Ph_{h,b-2D(h)}(Z_{h,b-D(h)}Z_{h,b}f)$.
Adding the three contributions in \eqref{eq:lap} and collecting the two radial
ones gives \eqref{eq:laplace-weights}.
\end{proof}

The radial group closes on a \emph{single} weight, $b-1$: the second-order
radial derivative loses two powers through $\partial_s^{2}$ and recovers one
through the Jacobian factor $s=qX$. The axial group does not: its degree is
\begin{equation}
  b-2D(h)=b-(1-2h),
  \label{eq:degrees}
\end{equation}
and the two degrees differ by
\begin{equation}
  (b-2D(h))-(b-1)=1-2D(h)=2h\ \in\ (0,1).
  \label{eq:deficit}
\end{equation}
Neither degree is the weight an isotropic Laplacian would carry in $q$, which is
$b-1$: that is the radial degree itself, and the axial degree equals it only at
$h=0$. For $q\in(0,1)$ the axial coefficient is the \emph{smaller} of the two by
the factor $q^{2h}$, so the axial term is the more damped of the two and the
radial term is the more singular near the apex. Which contribution dominates a given
expression also depends on the inner factors $Z_{h,b-D}Z_{h,b}f$; the weight
alone does not decide it, and a bound that treats $\nu\Delta$ as a single
homogeneous object of degree $b-1$ mis-assigns the axial piece.

\begin{corollary}[No admissible member is isotropic]
\label{cor:noiso}
Let $h\in(0,\tfrac12)$. Then $2D(h)=1-2h\in(0,1)$. Consequently the axial
degree $b-2D(h)$ differs from the radial degree $b-1$ for every admissible $h$:
the coincidence $2D(h)=1$ holds exactly at $h=0$, and $2D(h)=2$ at
$h=-\tfrac12$, both outside the admissible range of every declaration in
\file{SimilarityProfile.lean}, which all carry $0<h$ and $h<\tfrac12$.
\end{corollary}

\begin{proof}
$D(h)=\tfrac12-h$ gives $2D(h)=1-2h$; the bounds follow from $0<h<\tfrac12$.
The equations $1-2h=1$ and $1-2h=2$ are $h=0$ and $h=-\tfrac12$.
\end{proof}

\begin{remark}[What this does and does not say]
\label{rem:scoping}
Corollary~\ref{cor:noiso} is not a claim that the construction fails to be
self-similar; it is a statement about the \emph{operator}: on this family the
viscous term is a sum of three homogeneous pieces of two distinct degrees, and the
degree of the axial piece is $h$-dependent. Any estimate written as
``$\nu\Delta$ acts with weight $q^{b-2}$'' is, on this family, an assertion about
the radial part alone. Whether the accompanying manuscript makes that assertion
is a question about the manuscript's text; what we can certify is that the code's
own identities \eqref{eq:dss}--\eqref{eq:dzz} forbid the isotropic reading for
every $h$ in the range the code uses.
\end{remark}

\subsection{The container absorbs the mismatch}
\label{sec:container}

The anisotropy is not a defect of the coordinates; it is the geometry of the
cone doing its job. The inner domain is $\{\abs{\eta}<1\}$ with $X$ ranging over
a bounded interval for a compactly supported profile, so by \eqref{eq:inner} the
physical support satisfies
\[
  \dd z\ \text{scales as}\ q^{D(h)},\qquad \dd s\ \text{scales as}\ q ,
\]
and the physical measure element is
\begin{equation}
  \dd z\,\dd s\ =\ q^{\,1+D(h)}\,\frac{L(h,\eta)}{d(\eta)}\ \dd X\,\dd\eta ,
  \label{eq:measure}
\end{equation}
because $z=q^{D}\eta$ with $q\,d(\eta)=\tau$ gives
$\dd z=q^{D}(L/d)\,\dd\eta$ at fixed $\tau$, while $\dd s=q\,\dd X$. Hence
\begin{equation}
  \bigl\lVert \Ph_{h,b}f\bigr\rVert_{L^{2}(\dd z\,\dd s)}^{2}
  \ =\ q^{\,2b+1+D(h)}\int_{\abs{\eta}<1}\frac{L(h,\eta)}{d(\eta)}\,
  \abs{f(X,\eta)}^{2}\,\dd X\,\dd\eta ,
  \label{eq:l2}
\end{equation}
so the weight that makes the \emph{scaling} of the norm cancel is
\begin{equation}
  b_{*}(h):=-\frac{1+D(h)}{2}=-\frac{3-2h}{4}.
  \label{eq:bcrit}
\end{equation}
At $b=b_{*}$ the pullback is scale-invariant but not an isometry: the inner
integral still carries the weight $L/d$. On a profile supported in
$\abs{\eta}\le1-\varepsilon$ that weight is bounded between two positive
constants, since $L(h,\eta)=1-2h\eta^{2}\ge1-2h$ and
$d(\eta)=1-\eta^{2}\ge2\varepsilon-\varepsilon^{2}$; without that separation
from the cone boundary the equivalence fails as $\abs{\eta}\to1$.

At $h\to\tfrac12$ this is $-\tfrac12$ and at $h\to0$ it is $-\tfrac34$; the
critical weight moves with the cone's opening. The viscous losses sit at
$b_{*}-1$ (the radial group) and $b_{*}-2D(h)$ (the axial group), so the axial
term is the more damped of the two, and the radial term is correspondingly the
more singular.

\begin{remark}
Equations \eqref{eq:measure}--\eqref{eq:bcrit} are not stated in the repository;
they are immediate consequences of \eqref{eq:inner}, \eqref{eq:scaling} and
Definition~\ref{def:pullback}, and we label them as derived. They are included
because they make the geometric reading of Proposition~\ref{prop:three}
explicit: the generalized cone is a container whose axial opening $D(h)$
determines both which derivative loses least and which weight is $L^2$-critical.
An exponent ledger that assigns a single viscous exponent is, in this geometry,
assigning the wrong bookkeeping to one of the three terms.
\end{remark}

\subsection{Why the exponent ledger cannot see this}
\label{sec:ledger-cannot}

There is a precise reason this anisotropy survives the formalization, and it is
the same reason the constants do (\S\ref{sec:index-vs-term}). The weight $b$ in
Definition~\ref{def:pullback} is a \emph{real number argument} of a definition;
it occurs in the term $q^{b}$ via \ct{Real.rpow}, whose type is
$\R\to\R\to\R$ regardless of $b$. Changing $b$ to $b-2h$ does not change any
type, does not create a unification problem, and cannot be detected by the
kernel. Likewise $D(h)$ is a \emph{definition} (\ct{CoordinateAlgebra.D}), and
the identity $2D(h)=1-2h$ is proved by \ct{ring}; but no declaration asserts
$2D(h)\neq2$, because that is a statement about the geometry, not about the
algebra the file sets out to check.

So the ledger's arithmetic is exact where it operates---on rationals, by
\ct{linarith} and \ct{ring}---and silent where the geometry enters. The
anisotropy of \eqref{eq:laplace-weights} is visible in the code, in four
theorems with explicit weights; what is absent is any declaration that would
\emph{fail} if one treated those weights as equal. That absence is structural: to
make $b$ an index one would have to reflect it into a type, and the repository
does not.

% =====================================================================
\section{Three interface audits}
\label{sec:interfaces}

The findings of \S\ref{sec:constants} and \S\ref{sec:anisotropy} concern the
\emph{interior} of the analytic chain. A separate pass audited the three
\emph{interfaces} at which a construction of this kind normally pays a visible
price: the boundary at spatial infinity (where a pressure--velocity pairing is
integrated by parts), the temporal junction at the rest margin $t=3/8$ (where a
field switches on), and the divergence defect produced by cutting a
divergence-free field. In each case the audit found the same pattern: \textbf{what disappears is a
separate boundary, matching, or corrector estimate---not because the relevant
term is small, but because each is replaced by an identity.} Each interface is
closed by a germ equality, a case split, or a curl rather than by an inequality,
which is why none of the three contributes a constant, an exponent, or a
hypothesis to the ledger. The three are not free: \S\ref{sec:divergence-corrector}
records where the cutoff commutator is actually paid for.

\subsection{Pressure non-locality and the spatial decay boundary}
\label{sec:riesz-decay}

Pressure is the non-local solution of $\Delta p=\partial_i\partial_j(u_iu_j)$.
Any argument that integrates $\nabla p$ against a field must therefore either
assume decay at infinity,
$\lim_{R\to\infty}\int_{\partial B_R}p\,(w\cdot n)=0$, or restrict the
competitor class to fields of compact support. The audit asked which of the two
the repository does.

\textbf{Neither.} The declaration \ct{classical\_uniqueness\_on\_Icc}
(\file{NavierStokes/R3/WholeSpaceUniqueness.lean:30}) contains no boundary term
to evaluate. Its hypotheses are smoothness of $u,v,p,q$ on the slab $[0,T]$;
\emph{compact spatial support of the reference $u$ only} (\ct{hK},
\ct{hsupp}); divergence-freeness of both fields; equality of the two residuals;
equality of the data; and a uniform finite-energy bound on the competitor $v$
(\ct{hev}). There is no support, decay, growth, or derivative hypothesis on $v$,
and no bound, decay, or growth hypothesis on $q$.

Two mechanisms make the surface term vacuous. First, the pressure never appears
alone: it is paired with $\nabla\chi_R$,
\[
  \int_{\R^3}(p-q)\,\nabla\chi_R\cdot w ,
\]
and $\chi_R$ is compactly supported for every finite $R$, so no integral over
$\partial B_R$ is ever written. Second, the Riesz identity
\ct{actual\_flux\_eq\_canonicalCutoffFlux} rewrites that pairing as the canonical
Riesz functional of the tensor $u\otimes v$, after which
\ct{PressureFlux.exists\_uniform\_actual\_pressure\_flux\_bound} bounds it by a
constant that is fixed before $R$ is quantified. Every remaining $R$-dependence
in the comparison is explicit and vanishes:
\[
  \frac{C_0}{R^2},\qquad \frac{C_1}{R}B^{3/2},\qquad
  C_2\Bigl[\bigl(B^{1/2}+1\bigr)\Bigl(\frac{A}{R}+\frac1{R^2}\Bigr)
  +R^{-7/4}B^{3/4}\Bigr],
\]
(\file{NavierStokes/R3/ComparisonRateBound.lean:130}); the closure
\ct{eq\_of\_pressure\_flux\_bound} then drives $\norm{w(t)}_2^2\to0$ by
dominated convergence.

So the boundary term is neither assumed away nor bypassed by a qualitative
support predicate on the competitor: it is \emph{never formed}. Spatial decay at
infinity is replaced by two weaker and checkable inputs---finite energy of the
competitor, compact support of the reference---and the limit $R\to\infty$ is
taken on a bound, not on an integral identity. This is a stronger conclusion
than the one the audit task anticipated, and it is the reason no decay lemma
exists to be found: \file{WholeSpaceUniqueness.lean} contains no such statement
because the proof never needs one.

\subsection{Temporal interface regularity at the initial rest margin}
\label{sec:temporal-bump}

The second interface is the junction at $t=3/8$, where the constructed field
passes from identically zero to active. The audit asked whether smoothness there
is carried by a global \ct{ContDiff} statement or by a piecewise conditional that
delegates smoothness to sub-intervals, and whether the force's time derivatives
are evaluated explicitly at $t=3/8$.

\textbf{Neither.} There is no \ct{if-then-else}, and no derivative is evaluated
at a point. The junction is a scalar multiplier in time alone:
\ct{activatedVelocity}\,$u:=\lambda(t)\,u$ and
\ct{activatedPressure}\,$p:=\lambda(t)\,p$
(\file{NavierStokes/TimeLocalization.lean:27}--\file{:31}), together with
\ct{PositiveTimeForce.force}\,$f:=\theta(t)\,f$
(\file{NavierStokes/R3/PositiveTimeForce.lean:46}). Smoothness is inherited by
closure of $C^\infty$ under multiplication and composition:
\ct{activatedVelocity\_smooth} and \ct{force\_contDiff} are one-line
applications of \ct{ContDiffOn.smul} and \ct{ContDiff.smul}. The two branches
agree because the field itself vanishes on the early slab.
\ct{activatedVelocity\_zero\_early} and \ct{activatedPressure\_zero\_early} give
$u=p=0$ for $|t|\le3/8$; \ct{localized\_residual\_zero\_early} then gives a
vanishing residual there, and the equation across the junction is closed by a
two-branch case split on $t<3/8$
(\file{NavierStokes/R3/ActualCandidate.lean}): on the early side both residual
and force are zero, on the late side \ct{PositiveTimeForce.force\_eq} applies on
$[3/8,1]$.

Two clarifications correct the premise of the task. The declaration it named,
\file{NavierStokes/R3/CandidateBreakdown.lean}, contains no smoothness proof of
$f$ at all: its four theorems are the uniqueness transfer, the force-nonzero
lower bound, and an $L^2$ bound. Force smoothness is established elsewhere, as
\ct{compactForce\_smooth} (a smooth spatial outer cutoff times $f$,
\file{NavierStokes/R3CompactCandidate.lean:90}) followed by
\ct{PositiveTimeForce.force\_contDiff}. And the force's time support is not
asserted near the junction but bounded away from it:
\ct{force\_tsupport\_subset} gives
$\operatorname{tsupp}f\subseteq[1/16,21/16]\times K$, which is what supplies the
\ct{CompactPositiveTimeSupport} field.

The inheritance is only half of that obligation, and the other half is where
the sharpest external objection landed: the final assembly also requires the
residual's jets to converge locally uniformly at the singular point
(\ct{hlim}, \file{NavierStokes/CandidateFromLimits.lean:39}). That convergence
is not postulated. Vanishing joint jets at the singular spatial point, a common
schedule carrying residual flatness (\ct{exists\_schedule\_all\_jetRates}), and
locally uniform boundary limits (\ct{boundaryLimits\_locallyUniform}) assemble
it from the finite-stage estimates; the chain, and the disposition of
\ct{meanSolve} and \ct{highSolve}---Euler-only, not on this path---are E65--E66.
The estimates themselves are then not merely consumed: they are traced below
\ct{RawStageBounds} to a single wave-class fact about the actual source (E67,
E68), that fact is proved for the initial state and rebuilt at every cycle
(E69), and the assembly which rebuilds it is a structure literal with no
assumed field (E71).

Consequently no jet of $f$ is matched at $t=3/8$ by computation. The junction is
regular because both sides are eventually equal to the zero germ, and
$C^\infty$ is closed under that identification.

\subsection{Divergence-free cutoff corrections in the torus interface}
\label{sec:divergence-corrector}

Cutting a divergence-free field creates a defect:
$\nabla\cdot(\chi u)=u\cdot\nabla\chi\neq0$. The standard repair is a
corrector field, obtained by solving a divergence equation and subtracting a
gradient. The audit asked whether such a corrector exists, whether its spatial
derivatives consume powers of $q$ missing from the exponent ledger, and whether
its residual is added to $f$ or majorized by an existential constant.

\textbf{There is no corrector, and none is needed.} The construction cuts the
\emph{potential}, not the velocity:
\[
  \ct{cutPotential}\,A:=\chi\,A,\qquad
  \ct{cutVelocity}\,A:=\nabla\times(\chi A)
\]
(\file{NavierStokes/SpatialLocalization.lean:165}--\file{:172}), and a curl is
divergence-free identically, \ct{spatialDivergence\_spatialCurl\_on}
(\file{NavierStokes/SpatialLocalization.lean:260}). The periodized velocity is
the curl of the periodized potential,
\ct{periodicVelocity}\,$A:=\nabla\times\ct{periodicPotential}\,A$
(\file{:213}), with \ct{periodicVelocity\_divergence\_free} giving
$\nabla\cdot u=0$ exactly. The quantity $\nabla\cdot(\chi u)$, for which a
corrector would be required, therefore never appears in the development.

The price is paid elsewhere and is visible. The cutoff gradient is not removed;
it is \emph{retained inside the velocity}, by the product rule
\[
  \nabla\times(\chi A)=\chi\,\nabla\times A+\nabla\chi\times A .
\]
The second term is precisely the commutator $[\nabla,\chi]$ that a corrector
would have been introduced to cancel. Here it is part of the field, which is why
the cut velocity is exactly divergence-free while being only germ-equal to
$\chi\,\nabla\times A$: on the plateau, where $\chi\equiv1$, the commutator
vanishes, and \ct{periodicVelocity\_eventuallyEq}
(\file{NavierStokes/SpatialLocalization.lean:318}) upgrades the plateau identity
to equality of all jets. Off the plateau the commutator is carried by
$\nabla\chi\times A$, supported in the annulus
$\ct{supportCylinder}\setminus\ct{plateau}$
(\file{:72} and \file{:134}), and it enters the residual as an exact term of the
equation.

Both ledger questions therefore have vacuous answers, and that is the finding:
\file{ExponentLedger.lean} has no corrector entries because there is no
corrector, so no powers of $q$ can be missing from it; and the corrector's
residual is not absorbed into $f$ by a majorization, because there is no
corrector residual to absorb. The force remains, as established in
\S\ref{sec:twohalves}, the exact residual of the fields actually used, and the
equation holds by case split rather than by estimate.

% =====================================================================
\section{Packet time-width compression: derived, with one conditional link}
\label{sec:packets}

The findings of \S\ref{sec:constants} and \S\ref{sec:anisotropy} concern the
\emph{statement} of the theorem and the \emph{homogeneity class} of its
operator. Neither says anything about how the existence half is assembled:
\S\ref{sec:twohalves} takes a field, names its residual, and stops there,
leaving the provenance of the field and the smoothness of its residual to
the assembly. This section and the next close that
gap in the repository's own terms, on the Euler side of the development
(\file{Euler/}), where the packet family is built and where the top-level
singular datum is defined (audit log, E53).

\subsection{The object is a packet, not a pulse}

A scan keyed to the vocabulary of oscillatory pulses returns nothing: a
recursive search of \file{Euler/} and \file{NavierStokes/} for the strings
\ct{pulse} and \ct{oscillat} is empty, and no file name in either tree contains
either word. The construction
is named \emph{packet}. The subtree is large---\file{Euler/} holds $1{,}839$
files against $816$ in \file{NavierStokes/}, of which $110$ are
\file{Euler/Packet*.lean} mentioning compact support or a support
inclusion---so the naming question is not cosmetic. A scan for ``pulse'' reports
the absence of a name, not the absence of a mechanism, and the distinction is
the difference between a negative audit result and a false one. What is proved
here is a construction, not a continuation: the packets' time windows halve, and
the datum they define is the limit. Nothing in this section asserts smooth
continuation of an evolving solution through $t=1$; the smoothness proved is
that of the limiting initial datum (E57).

\subsection{The halving is a theorem about one geometric sequence}

The object that shrinks is \ct{timeWidth} (\file{Euler/PacketSourceScaleSequence.lean:42}):
\begin{equation}
  \ct{timeWidth}\ J\ X\ n =
  \frac{3\cdot\ct{scaleSequence}\ J\ X\ (n+1)\cdot\ct{scaleSequence}\ J\ X\ n}
       {\sqrt{\ct{previousShear}\ J\ X\ n}}\,,
\end{equation}
where \ct{scaleSequence} is geometric with $\ct{scaleSequence}\ J\ X\ 0 = X$ and
$\ct{scaleSequence}\ J\ X\ (n+1) = ((J+n:\mathbb{N}):\R)^{2}\cdot\ct{scaleSequence}\ J\ X\ n$
(\file{Euler/PacketSourceScaleChoice.lean:226--232}), and \ct{previousShear} is
exponential in \ct{scaleSequence} at exponent $7$. The numerator is a product of
consecutive terms; the denominator is super-exponential. Nothing asymptotic is
required to see the ratio fall.

\ct{actualWidths\_contract} (\file{Euler/PacketSourceScaleGuards.lean:69})
concludes
\[
  \ct{timeWidth}\ J\ X\ (n+1)\;\le\;\ct{timeWidth}\ J\ X\ n\,/\,2
\]
from $3\le J$, $8\le X$, an \ct{ActualBounds} witness, and $\delta\le 1/2$. The
route is \ct{actualTimeRatio\_le} (\file{:16}), which bounds the width ratio by
\ct{sourceTimeRatio}---one component of the \ct{ActualBounds} series---after
which \ct{SmallSeries.term\_le} yields ratio $\le\delta\le 1/2$ and
\ct{div\_le\_iff\_0} converts that into halving. A finite series total, then
arithmetic: no limit, no compactness, no asymptotic expansion.

\subsection{One link is a hypothesis, and the join is made elsewhere}

\ct{horizonTime\_antitone} (\file{Euler/PacketNestedHorizons.lean:90}) carries
\[
  (hw : \forall n,\ \ct{timeWidth}\ J\ X\ (n+1) \le \ct{timeWidth}\ J\ X\ n/2)
\]
as a binder, and that file does not import the guards file. Read alone, the
nesting theorem is therefore conditional, and citing it as an unconditional
compression result overstates by one import. The join is made through the guard
structure: \ct{StageGuards.next\_width}
(\file{Euler/PacketSourceScaleGuards.lean:129}) is a structure field,
\ct{stage\_guards} fills it with \ct{actualWidths\_contract}, and
\ct{exists\_guarded\_sequence} produces $J$, $X_0$, $\delta$ uniformly. The
compression closes through \ct{StageGuards}, not through
\file{PacketNestedHorizons.lean} alone. Section~\ref{sec:root} shows where that
field is consumed.

The arithmetic the nesting consumes is \ct{nested\_horizon\_of\_width\_ratio}
(\file{Euler/EulerProof.lean:19529}): from $\Delta\le 2W/3$ and
$W_{\mathrm{next}}\le W/2$ it yields $t+\Delta+2W_{\mathrm{next}}\le t+2W$, with
$\Delta$ bounded by \ct{stepLength\_bounds}
(\file{Euler/PacketNestedHorizons.lean:82}). Those constants give a strict drop
that the library does not state: since $2/3+1-2=-1/3$,
\[
  \ct{horizonTime}\ (n+1)\;\le\;\ct{horizonTime}\ n-\ct{timeWidth}\ n/3\,.
\]
The library records only $\le$, which is all the nesting argument needs.

\subsection{Grade is the derivative order}

Whether ``the time-supports compress as the derivative index rises'' is true of
\emph{this} construction depends on whether the grade index $n$ is the jet order
rather than a bare stage counter. It is, in the operative sense. \ct{velocityJet}
is assembled from \ct{slicedJet} applied to \ct{(a i).high}, \ct{(a i).mean} and
\ct{(a (i-1)).corrector} (\file{Euler/PacketProfileRecursion.lean});
\ct{knownJets}\ O\ p\ a is \ct{history}\ p of that jet family, so grade $p$ is
exactly the order-$p$ jet being cancelled; \ct{profiles\_step} for $p\ge 2$
depends only on the strict prefix, which is the well-founded recursion; and
\ct{tail\_grade\_bound} (\file{Euler/PacketProfileTailGrade.lean:29}) carries
$N+1\le n$, i.e.\ tail grades indexed beyond the required order $N$.

The literal theorem is nevertheless about consecutive \emph{grades} $n$ and
$n+1$. \ct{timeWidth} has no separate derivative-order parameter; the two
indices coincide by construction, not by definition. A reader who takes the
sentence as a statement about $\partial^{k}$ for arbitrary $k$ is reading one
step further than the code does.

\subsection{Three things the compression is not}

\begin{enumerate}
\item \textbf{Not a dimensional law.} Both the grading and the shrinking come
from the same \ct{scaleSequence}; the ratio is controlled by a summable cost
series, not by differentiating a scaling ansatz and reading off exponents. It
is therefore not evidence for or against the isotropy question of
\S\ref{sec:anisotropy}.
\item \textbf{Not the spatial support.} The spatial radius is \ct{supportScale},
exponential at exponent $7/2$, and \ct{scale\_support} and \ct{high\_scaled\_support}
(\file{Euler/PacketInitialSupport.lean:24}, \file{:48}) dilate a fixed ball by
$\ell$. Time width and spatial radius are distinct objects with distinct
exponents, $7$ and $7/2$; a claim about one is not a claim about the other.
\item \textbf{Not collapse to a point.} \ct{activationTime} is strictly
increasing while \ct{horizonTime} is antitone, so the intervals nest from both
ends---but \ct{common\_positive\_interval} proves
$\ct{baseHorizon}\ J\ X/12 < \ct{horizonTime}\ J\ X\ n$, so the intersection
has positive length. Nesting is used to keep one common interval available, not
to squeeze it shut.
\end{enumerate}

% =====================================================================
\section{The \texorpdfstring{\ct{Classical.choice}}{Classical.choice} that selects the scale regime}
\label{sec:root}

The compression theorems of \S\ref{sec:packets} are polymorphic in a scale
record \ct{S : Scales c B} and consume its \ct{ActualBounds} field. That raises a
question the local proofs cannot answer: does such a record actually reach the
top-level theorem, or is the field merely asserted somewhere adjacent to it? The
answer is that the chain closes through a single opaque selection of the scale
regime (audit log, E54). The lifespan is a second \ct{choose}, downstream of
that selection (\file{Euler/PacketFiniteLifespan.lean:54}), not a competing root.

\subsection{The five links}

\begin{enumerate}
\item \textbf{Existence.} \ct{exists\_scales}
(\file{Euler/PacketInductionScales.lean:123}):
\[
  \ct{theorem exists\_scales}\ (c\ B:\R)\ (hc: 0\le c) : \ct{Nonempty}\ (\ct{Scales}\ c\ B)
\]
Its proof takes \ct{exists\_base\_power}, then
\ct{exists\_common\_guards}, then assembles the record with
\ct{actual := hdata.1}. An \ct{ActualBounds} witness is \emph{produced}, for some
$J,D,X,\delta$, from nothing but $0\le c$.
\item \textbf{Selection.} \ct{constructionScales}
(\file{Euler/PacketInfiniteConstruction.lean:68}):
\begin{lstlisting}
noncomputable def constructionScales : ConstructionScales :=
  Classical.choice (exists_scales (requiredExponent : ℝ)
    (commonThreshold gradientConstant hessianConstant) (Nat.cast_nonneg _))
\end{lstlisting}
\item \textbf{Packets.} \ct{packets}\ n = \ct{stages}\ \ct{constructionScales}\ \ct{le\_rfl}\ \ct{le\_rfl}\ n
(\file{:72}). The two reflexivity arguments discharge the \ct{stages} hypotheses
$\ct{requiredExponent}\le q$ and $\ct{commonThreshold}\le B$, so neither is
carried downstream.
\item \textbf{Datum and lifespan.} \ct{initialDatum} =
\ct{Stage.initialDataLimit}\ \ct{packets} (\file{Euler/PacketFiniteLifespan.lean:18});
\ct{lifespan} = \ct{initialDatum\_finite\_lifespan.choose} (\file{:54}), with
\ct{lifespan\_le\_base} bounding its duration by
\ct{baseHorizon}\ \ct{constructionScales.J}\ \ct{constructionScales.X}.
\item \textbf{Top level.} \ct{initialDatum\_singularity}
(\file{Euler/EulerSingularity.lean:115}) and
\ct{exists\_compact\_smooth\_euler\_singularity} (\file{:133}).
\file{Euler.lean} is one line, \ct{import Euler.EulerSingularity}, so that file
is the root.
\end{enumerate}

\subsection{Where the compression is spent}

Each \ct{Stage} carries \ct{source\_stage : StageGuards}\ \dots
(\file{Euler/PacketInductionStage.lean:123}), obtained from \ct{Scales.stage},
which passes \ct{S.actual} into \ct{stage\_guards\_at} $\to$ \ct{stage\_guards}
$\to$ \ct{actualWidths\_contract}. That is the only place
\ct{next\_width} is proved. \ct{PacketStageRestriction.nextHorizon\_lt}
(\file{Euler/PacketStageRestriction.lean:49}) consumes it
(\ct{have hwidth := P.source\_stage.next\_width}) to place the next horizon
strictly below the parent's terminal time; without the halving that inequality
does not close. The remaining containment comes from
\ct{Stage.horizon\_eq}, which reads
$\ct{parent.T}=\ct{time}+2\cdot\ct{timeWidth}\ n$; from \ct{Stage.horizon\_le},
which reads $\ct{parent.T}\le\ct{baseHorizon}$; and from \ct{Scales.time\_small},
which reads $\ct{baseHorizon}\le 1$. Together they keep every packet inside
$[0,1]$.

\subsection{Instrumental, not constitutive}

The exact statement of \ct{initialDatum\_singularity} is worth quoting, because
it settles what the choice is \emph{for}. The theorem asserts smoothness
(\ct{ContDiff} at order $\infty$), compact support, non-vanishing, vanishing divergence at
every point, a positive duration bounded by $1$, the equivalence ``smooth
solution at $T$ $\iff$ $0<T<\ct{lifespan.duration}$'', a $C^{1}$ limsup equal
to $\top$, and a maximal-vorticity integral equal to $\top$. Nothing in that
list mentions \ct{timeWidth}, \ct{ActualBounds}, or any scale parameter. The
halving is used instrumentally, to build the \ct{packets} sequence whose limit
the datum is; it is not a conjunct of the conclusion. So the chain closes in the
construction and not in the statement, which is precisely why a reader who opens
only \file{EulerSingularity.lean} cannot see it, and why the trace is worth
recording.

\subsection{Logical closure, numerical opacity}

Because \ct{constructionScales} is a \ct{Classical.choice}, the parameters
$J$, $D$, $X$, $\delta$ have no computable content, and every downstream
quantity---\ct{timeWidth}, \ct{scaleSequence}, \ct{baseHorizon}, every
frequency---is defined from that opaque witness. The halving theorem is
nevertheless \emph{about the actual object}: every lemma is polymorphic in
\ct{S : Scales c B} and uses only fields, so nothing is assumed of a
hypothetical record. What is lost is recoverability, not correctness: no
explicit bound on the first width, on $J$, or on \ct{baseHorizon} follows from
the statement.

This is the barrier of \S\ref{sec:constants} at a coarser grain. There the
opacity sits behind individual constants; here it covers the entire scale
regime---chosen once at the root and inherited by every packet. The difference
is that this opacity is structural rather than incidental: the choice
\emph{is} the construction, and no rewriting of the proof would expose the
witness without also exposing which witness was chosen.

Two incidental checks are worth recording. First, the record's
\ct{delta\_small} ($\delta\le 1/16$) and \ct{delta\_geometry}
($10^{6}\cdot\ct{geometryConstant}\cdot\delta\le 1$) are stronger than the
$\delta\le 1/2$ that \ct{actualWidths\_contract} requires, so the halving has
slack to spare. Second, \ct{ActualBounds} has further producers, at
\file{Euler/PacketSourceScaleActual.lean:230} and
\file{Euler/PacketSourceScaleFromCosts.lean:70}; both are theorems, and no path
enters \ct{ActualBounds} as an axiom or a bare hypothesis.

Both questions this trace originally left open are now closed. The import path
is four hops---\file{EulerSingularity} $\to$ \file{EulerC1Limsup} $\to$
\file{EulerC1Breakdown} $\to$ \file{EulerFiniteLifespan} $\to$
\file{PacketFiniteLifespan}---and \ct{PacketFiniteLifespan} has exactly one
direct importer inside the root's reachable set, so the final edge is forced
(E55). The chosen $D$ does not constrain the derivative order: it enters
\ct{previousFrequency} only at index zero, as $X^{D}$, and every later term is
$D$-free. Its one degrading effect is the index-zero instance of
\ct{neighborError}, which carries $X^{cD}$ and is absorbed by enlarging $X$;
there is no upper bound on $D$ anywhere in \ct{ActualBounds}, so no ceiling on
the order follows (E56, corrected in E59~\S3). The width argument never reads
it.

% =====================================================================
\section{Conclusion}
\label{sec:conclusion}

\subsection{The verdict}

The logical integrity of \texttt{openai/NavierStokesAndEuler} is robust. There
are no declared axioms, no \ct{native\_decide}, no \ct{unsafe}, no \ct{partial}
definitions, no \ct{implemented\_by} hooks, and no reachable \ct{sorry}; the
exponent ledger is exact rational arithmetic; and the comparison theorem that
supplies the nonexistence conjunct is a genuine uniqueness result whose
hypotheses are nearly empty---no pressure norm, no decay of $p-q$ at infinity, no
support, growth, or derivative condition on the competitor, with the Gronwall
constant uniform in the localization radius by binder order and closure by an
interior-derivative argument that uses $E(0)=0$ at a single point. That is not a
proof that has been padded into shape.

But its success is an achievement of \emph{statement design}, not of fluid
physics, and the structural findings of this paper delimit what the artifact
certifies:

\begin{enumerate}
\item \textbf{The Non-Evaluated Constant Barrier.} Every quantitative constant
is $\exists$-bound or \ct{choose}n, so the kernel verifies existence and never
size or dimension (\S\ref{sec:constants}). The headline theorem is
constant-free, so it is not threatened; but no effective bound is extracted, and
no machine check of dimensional consistency is possible.
\item \textbf{Anisotropic viscous scaling.} The Laplacian on the generalized
cone is a sum of two homogeneous pieces with degrees $b-1$ (radial) and
$b-2D(h)$ (axial), where $2D(h)=1-2h\in(0,1)$, differing by $2h$; the axial
degree reaches the isotropic weight $b-1$ only at $h=0$, so no admissible member is isotropic, and the axial
term is the more damped of the two
(\S\ref{sec:anisotropy}). The container $|z|<q^{D(h)}$ is what absorbs the
mismatch, through the measure factor $q^{1+D(h)}$ and the $h$-dependent
critical weight $b_{*}(h)=-(3-2h)/4$.
\end{enumerate}

Both are instances of one principle: \emph{the kernel audits arithmetic
reflected into types and is silent about arithmetic that is not.} Exponents that
reach an index are checked by unification; constants in existential binders and
real exponents inside definitions are not. A reader who takes ``machine-checked''
to mean ``numerically and dimensionally reliable'' is therefore misreading the
artifact---not because the artifact is defective, but because that is not what
type checking is.

A third observation is of a different kind. The three interfaces audited in
\S\ref{sec:interfaces} require no separate boundary, matching, or corrector
estimate: the pressure pairing has no boundary
term at infinity because the test is compactly supported and the limit is taken
on a vanishing bound; the temporal junction is a smooth scalar multiplier closed
by a vanishing germ rather than by a jet match; and the divergence defect of
cutting is avoided by curling the cut potential, which retains the commutator
$\nabla\chi\times A$ inside the velocity instead of correcting it away. None of
the three appears as a constant, an exponent, or a hypothesis. The interfaces of
this construction are exact identities, and exactness is why the
constant barrier of \S\ref{sec:constants} never has to be crossed there.

The construction side carries its own ledger, and Sections~\ref{sec:packets} and
\ref{sec:root} trace it. Packet time widths halve by theorem rather than by
hypothesis---\ct{actualWidths\_contract} derives
$\ct{timeWidth}\,(n+1)\le \ct{timeWidth}\,n/2$ from a finite series total---with
one conditional link, \ct{horizonTime\_antitone}, joined not inside its own file
but through the \ct{StageGuards.next\_width} field. And the entire scale regime
$J$, $D$, $X$, $\delta$ enters the development through a single
\ct{Classical.choice} at \file{Euler/PacketInfiniteConstruction.lean:68}, from
which the packet family, the limiting datum and its lifespan are all built.
That root is why the final theorem can be constant-free in its statement while
being opaque in its construction: the compression is spent to build the object,
and the object's definition is the only place the witness is ever named.

\subsection{What the theorem says, in one sentence}

The theorem proves that a smooth compactly supported force can drive a fluid from
rest to unbounded speed with no global finite-energy solution. That is the data
and the conclusion of alternative (C) of the Clay statement: zero is an
admissible initial datum, and a smooth compactly supported force is stronger
than Clay's rapid-decay requirement. If the proof is correct, it resolves (C).
It does not establish the unforced alternatives (A) or (B), and the repository's
own \ct{force\_nonzero\_of\_no\_global\_solution} is the formal witness that the
result cannot be restated with $f\equiv0$---which is exactly why it bears on (C)
and not on (A). Whether the proof is correct, and how CMI evaluates it, are
separate questions that this audit does not decide.

\subsection{Scope and limitations}
\label{sec:limits}

The findings above are by inspection of proof terms, binder sequences and
definitions. Kernel acceptance was tested separately and afterwards: the pinned
toolchain (\ct{leanprover/lean4:v4.34.0-rc2}) installs on the audit host, but its
\ct{lake} aborts with \ct{SIGTRAP} in \ct{libmalloc} at worker-thread teardown
(reproduced 3/3), so the build ran in a Linux virtual machine (Ubuntu~26.04,
8 cores, 12\,GiB). \ct{lake build NavierStokes} completed 9,580 jobs with exit~0
in 20\,min\,17\,s, producing 817 modules; \ct{\#print axioms} on
\ct{breakdownStatement}, \ct{theorem\_1\_1},
\ct{theorem\_1\_1\_with\_initial\_rest} and
\ct{theorem\_1\_1\_with\_dissipation} reports only \ct{propext},
\ct{Classical.choice} and \ct{Quot.sound}. A later run from a pristine clone on a
GitHub Actions runner repeated the build (9,580 jobs, exit~0, 28\,min\,35\,s) and
printed five further footprints, the periodic theorem and its adapters among them
(E73). Comparator, the one external checker this audit ran, passes on the statement (E75). This reproduces acceptance for the
whole-space library; it does not re-run the other eight external checkers, and it
verifies \emph{terms}, not \emph{statements}. The statement-level caveats (zero
datum, force-as-residual, similarity in $\nu$) are exactly what kernel acceptance
cannot test. The derived identities
\eqref{eq:measure}--\eqref{eq:bcrit} are ours, not the repository's, and are
labelled as such. One entry of the audit log (E36, on
\ct{SchwartzCompactApproximation.continuous\_zero\_of\_compactSupport}) is
recorded as provisional: the call site and the approximation estimates were
verified, the exact binder was not opened. Finally, the claim in
\S\ref{sec:index-vs-term} that every use of a constant is monotone in that
constant is a meta-argument about the scripts and is not machine-checked.

\paragraph{The paper against the tree.} The repository links the prose account
it formalizes, and \file{formalization.yaml} names that account as its source,
aligning Theorem~1.1 and Corollary~10.6 with named declarations. Until E89 this
audit read the paper only where the statement comparison forced it, so the
question whether the argument formalized is the argument published had no
answer. E89 begins it: every numbered section of the 166-page Navier--Stokes
paper is mapped to the modules that realise it, and each correspondence is
labelled by its basis---\emph{cited}, when a Lean docstring names the section;
\emph{object}, when the module defines the object the section is about;
\emph{family}, when the modules cover the same subject and the correspondence is
unchecked. Five paper sections are cited in the source at all (\S4.4, \S4.6,
\S8.1, \S8.2, \S8.3, \S10.2), and eight of the thirteen rows rest on
\emph{object} or \emph{family}. The map therefore settles no lemma. It does not
show that each numbered result of the paper has a declaration, nor that no paper
lemma is missing from the formalization. That is the remaining work, recorded in
ile{ns-audit/CROSSWALK.md}. E91 put the paper into the tree as text and corrected
its count: 73 numbered results, not 57 (E135 corrected the kind breakdown that produced 74). Twenty-eight are now placed at statement level; one is partial, and forty-four are unplaced. The partial is Lemma~9.7, partial on statement shape rather than on content (E133). E151 places Lemma~8.8: the paper's $R_g$ is \ct{actualRadialError}, defined as the new radial source minus the old minus $(2V/R)\Delta v$; (8.26) is derived from the averaged equation by \ct{gr\_change}, (8.27) is the exact identity \ct{defects\_after\_solved\_rows} before any estimate is applied, and $H+9/10-2\kappa$ is $\alpha+0.9-2\kappa_s$ under the paper's own $9/10 \le H$. E152 places Proposition~8.3: \ct{meanPressure\_radial\_residual\_global} is (8.13) with both the bump mass and the cutoff alias displayed, \ct{meanPressure\_bar\_hasDerivAt} is the bar derivative derived from the integral representation, and the three stream classes are \ct{meanClass\_streamPotential}, \ct{meanClass\_streamGamma} and \ct{meanClass\_scaledRadialUpdate} at $\alpha$, $\alpha$ and $\alpha+1$, the last because the $\varepsilon$ sits inside the direction the derivative is taken along, exactly where the paper's $Q^{1/2-D}=\varepsilon$ puts it. E149 places Proposition~9.3: the band-independent harmonic set is \ct{modes residualBand}, one \ct{Finset} with the band quantified inside the statement and the set updated once per stage, and the base error's minus in (9.6) is on the opposite side of the equality rather than dropped. Proposition~6.6's last two clauses closed in E148 --- the pre-cutoff rectangle is the hypothesis form of \texttt{MeanLocalCoefficient.coefficient\_mul}, and the scaled normal row is \texttt{meanClass\_chartPrefactor\_all} --- and E146 and E147 placed Proposition~8.1 and Lemma~8.2. Lemma~9.2 left this column in E139, and §8 has none left at all: E140 and E141 closed Lemmas~8.7 and~8.6, and E145 closed Corollary~8.5, whose two radial divergence identities are declared at the paper's own $2/R$ and $1/R$ connection weights. E94 places Proposition~9.1, the row two earlier passes recorded as absent: the residual's class $W_{\alpha+1/2-3\kappa}$ and the pressure's $W_{\alpha+1/2}$ are one conjunction in \texttt{LocalizedWaveBounds}, stated for the object the tree calls \texttt{globalGood} and the paper calls the normalised linear residual. E92 places Proposition~9.9 whole --- the summed local field, with
divergence-freeness as a field of its own statement --- and finishes Theorem~4.6's
part~(vi), the two reserved intervals, in \texttt{ReservedPatches}: disjoint log
windows, $U=0$, and $E = c(1+\eta^2)^{-1}X^{-1/2-\lambda}$ on each. Part~(iv) was recorded by E92 as half placed; E93 closed it and corrected that
reading. The weight $\zeta$, its flatness at both edges and the $\delta^{-m_\alpha}$
loss are the tree's own \texttt{ActiveAnnulusWeight.weight} and \texttt{edgeDistance},
read at $y=\log X$. The uniform lower bound $|T_0| \ge c\zeta$ is
\texttt{LeadingStressWeights.stress\_lower\_bound}, with the same flat $\zeta$ and no
$\alpha$ in the statement. Its constant is not the cone margin --- a cone margin gives
only nonvanishing, which is \texttt{stress\_ne\_zero} --- but the minimum of two
edge-collar constants and a compactness margin over the closed annulus,
\texttt{exists\_global\_lower\_bound}. It is existential, as the paper asks, and no
numeral is exposed. Theorem~4.6 is therefore placed in all six parts.
Propositions~5.3, 9.5 and 9.6 were partial through E95 and are placed as of E97.
The arithmetic was never in doubt: the truncation exponent $2h(N+1)$ is
\texttt{slowOrder}, and $B_0 = 7/10$, $C^*_0 = 6/5$ with $1/10$ increments are
\texttt{residualWave} and \texttt{residualMean}. What E97 matched is the analytic
half clause by clause. Proposition~5.3's bound is
\texttt{truncationResidual\_rate}, whose exponent $2Nh-2-M$ is the paper's
$2h(N+1)-K_m$ at $K_m=m+2+2h$ --- exact in the derivative order and independent
of $N$ --- and its divergence clause is
\texttt{divergence\_slowVelocity\_eq\_zero}, stated for the actual Euclidean
divergence of the finite sum rather than in the similarity chart. Proposition~9.5's
mean class is \texttt{MeanResidualBounds}, which states $1+\sigma$ for every
admissible $\sigma$ and evaluates at $\sigma_0=1/5$ to the paper's $6/5$;
Proposition~9.6's cycle output is \texttt{gaugeTemporalStage\_next\_mean} at
$\sigma+1/10$, with the wave side at \texttt{gain} $h(j+1)$. One qualification is
carried rather than dropped: Proposition~9.6's four ordered corrections are proved
as four declarations, not as one composed step theorem, so placing it assumes the
composition.
The remaining gap is therefore breadth, not a single clause: forty-four numbered results are still unread at statement level, and the partial row is where the algebra is stated and the bounds are assembled --- the place a defect would sit, because the kernel never weighs the bounds. A kernel check cannot compare an argument against the argument it formalizes.

\paragraph{The §8 rows, and what a name costs.} E131 and E132 recorded Lemma~8.7 and
Lemma~8.6 as partial on specific negatives: no block inverse, no uniqueness declaration, no
seminorm-to-$C^m$ conversion. Each negative is a declaration one module away. The five-row inverses
are \ct{LocalizedMomentRepair.matrix} with \ct{det \neq 0} and \ct{coefficients} defined as the
inverse; the $S_\alpha \to M_\alpha$ clause is \ct{MeanRankUpdate.rank\_update\_meanClass};
uniqueness of the temporal inverse is \ct{TemporalMeanUpdate.timeDerivative\_unique}; and the
conversion is \ct{SmoothFourierData.coefficient\_seminorm\_bound} at $p+4$, composed into
\ct{ParametricTorusInverse.mixedJet\_inverse\_bound} with an explicit constant. One shape
difference is recorded rather than smoothed: the composite loss is $m+5$ where the paper prints
$m+4$, the extra order being composition order and recoverable from the tree's
$\text{weight}^{-4}$ summability but not printed. That is a limit on the printed index, not on the
estimate. E144 then re-read Proposition~9.3 against the periodized assembly: the decomposition, the
supports, the smooth zero-extension layer and (9.6) are declarations, and the row stayed partial on
exactly two clauses --- a single finite harmonic set carried across bands, and the sign of the
base-flat error. E149 closed both: the set is \ct{modes residualBand}, one \ct{Finset} quantified
outside the band, and the base error's minus is moved across the equality rather than lost.

\subsection{Implications beyond this repository}

Large formalizations of quantitative PDE estimates are becoming feasible, and
this artifact is a preview of their characteristic failure mode. The kernel will
faithfully audit the combinatorial and exponent-theoretic skeleton, which is
where human error is most common and most damaging. It will not audit the two
things a numerical analyst cares most about: the size of constants and the
homogeneity class of the operator. If such developments are to support claims of
the form ``the constant is explicit'' or ``the estimate is dimensionally
correct'', those claims need to be reflected into the types---as explicit
numerals, as rational bounds, or as homogeneity statements stated as theorems
about the pullback---because otherwise no amount of compilation will bear on
them.

% =====================================================================
\appendix
\section{Artifact index}

All paths relative to commit \texttt{f9e8bc5}. ``E'' references are entries of
\file{ns-audit/AUDIT\_LOG.md}.

\begin{center}
\small
\begin{tabular}{@{}llll@{}}
\toprule
Path & Declaration & Line & Used in \\
\midrule
\file{R3/ProblemStatement.lean} & \ct{breakdownStatement} & 150 & \S1.2 \\
\file{R3/ProblemStatement.lean} & \ct{CandidateProperties} & 92 & \S1.2 \\
\file{R3/Theorem.lean} & \ct{theorem\_1\_1\_with\_initial\_rest} & 26 & \S1.2 \\
\file{R3/Theorem.lean} & \ct{theorem\_1\_1} & 46 & \S1.3, Prop.~\ref{prop:constantfree} \\
\file{R3/CandidateBreakdown.lean} & \ct{force\_nonzero\_of\_no\_global\_solution} & 53 & \S1.3, \S5.2 \\
\file{R3/WholeSpaceUniqueness.lean} & \ct{classical\_uniqueness\_on\_Icc} & 30 & Rem.~\ref{rem:notfree}; used at 60, 64 \\
\file{R3/WeightedSobolev.lean} & \ct{sobolevConstant} & 25 & Ex.~\ref{ex:sobolev} \\
\file{R3/WeightedSobolev.lean} & \ct{lpNorm\_six\_le} & 42 & Ex.~\ref{ex:sobolev} \\
\file{R3/WeightedSobolev.lean} & \ct{weightedSobolevConstant} & 210 & \S3.1 \\
\file{R3/LocalizedFluxEstimates.lean} & \ct{exists\_weight\_second\_derivative\_bound} & 176 & Ex.~\ref{ex:choose} \\
\file{R3/LocalizedFluxEstimates.lean} & \ct{weightSecondDerivativeConstant} & 184 & Ex.~\ref{ex:choose} \\
\file{R3/PressureFunctionals.lean} & \ct{averagedPressureDifferenceValue\_bound} & 169 & Ex.~\ref{ex:pressure} \\
\file{R3/PressureFunctionals.lean} & \ct{averagedPressureDifference\_bound} & 279 & Ex.~\ref{ex:pressure}, E27 \\
\file{ExponentLedger.lean} & \ct{particularGain}, \ct{particular\_gain\_eq} & --- & \S3.5 \\
\file{ExponentLedger.lean} & \ct{manuscript\_kappa\_admissible} & --- & \S3.5 \\
\file{CoordinateAlgebra.lean} & \ct{D}, \ct{d}, \ct{L} & 19--21 & \S4.1 \\
\file{CoordinateAlgebra.lean} & \ct{timeCoeff}, \ct{axialCoeff} & 180--185 & Lem.~\ref{lem:pullback} \\
\file{SimilarityCoordinates.lean} & \ct{coordinateQ}, \ct{coordinateEta}, \ct{coordinateX} & 124, 340, 500 & \S4.1 \\
\file{SimilarityHomogeneity.lean} & \ct{coordinateQ\_scale\_h} et seq. & --- & \S4.1 \\
\file{SimilarityProfile.lean} & \ct{pullback} & 45 & Def.~\ref{def:pullback} \\
\file{SimilarityProfile.lean} & \ct{partialS\_pullback}, \ct{partialZ\_pullback} & 224, 233 & Lem.~\ref{lem:pullback} \\
\file{SimilarityProfile.lean} & \ct{partialS\_partialS\_pullback} & 279 & \eqref{eq:dss} \\
\file{SimilarityProfile.lean} & \ct{partialZ\_partialZ\_pullback} & 296 & \eqref{eq:dzz} \\
\file{PeriodicLocalization.lean} & \ct{innerCube} & 213 & \S2 \\
\file{PeriodicResidualLimits.lean} & \ct{periodicResidual\_jets\_eq\_cut} & 138 & \S2 \\
\file{R3/ComparisonRateBound.lean} & \ct{exists\_uniform\_rate\_bound} & 127 & \S1.3, \S5.1 \\
\file{SpatialLocalization.lean} & \ct{cutPotential}, \ct{cutVelocity} & 165--172 & \S5.3 \\
\file{SpatialLocalization.lean} & \ct{periodicPotential}, \ct{periodicVelocity} & 210--216 & \S5.3 \\
\file{SpatialLocalization.lean} & \ct{spatialDivergence\_spatialCurl\_on} & 260 & \S5.3 \\
\file{SpatialLocalization.lean} & \ct{periodicVelocity\_eventuallyEq} & 318 & \S5.3 \\
\file{TimeLocalization.lean} & \ct{activatedVelocity}, \ct{activatedPressure} & 27--31 & \S5.2 \\
\file{R3/PositiveTimeForce.lean} & \ct{force}, \ct{force\_eq}, \ct{force\_contDiff} & 46--55 & \S5.2 \\
\file{R3CompactCandidate.lean} & \ct{compactForce\_smooth} & 90 & \S5.2 \\
\bottomrule
\end{tabular}
\end{center}

\section{Non-claims}

For the avoidance of doubt, this paper does not claim, and its evidence does not
support, any of the following.

\begin{enumerate}
\item That the repository is unsound, contains a hidden axiom, or proves a false
statement. The verifications reported here point the other way.
\item That \Lean{} or Mathlib is responsible for the statement-design choices.
The same theorem, written in prose, would admit the same objections; and the
kernel is the component that behaved exactly as specified.
\item That the comparison theorem is weak or circular. It is, in our assessment,
the strongest part of the development, and its hypotheses are close to minimal.
\item That the anisotropy of \S\ref{sec:anisotropy} is an error in the code. The
code's identities are correct; the finding is that an isotropic reading of the
viscous term is incompatible with them.
\item That the whole repository is $641{,}332$ lines of comparison argument. That
figure is the whole tree; the comparison chain is a small fraction of it, and we
do not attribute the total to any single module.
\end{enumerate}

% =====================================================================
\begin{thebibliography}{99}

\bibitem{repo}
OpenAI, \emph{NavierStokesAndEuler}, GitHub repository, commit
\texttt{f9e8bc5} (code, formalization manifest, and accompanying manuscript).

\bibitem{audit}
A.~El Sayad and John, \emph{Forensic audit log for
\texttt{openai/NavierStokesAndEuler}}, \file{ns-audit/AUDIT\_LOG.md}, entries
E0--E97 (companion document).

\bibitem{lean}
L.~de~Moura and S.~Ullrich, \emph{The \Lean{}~4 theorem prover and programming
language}, Proc.\ CADE-28, Lecture Notes in Comput.\ Sci.\ \textbf{12699}
(2021), 625--635.

\bibitem{mathlib}
The Mathlib Community, \emph{The \Lean{} mathematical library}, Proc.\ CPP
2020, ACM, 2020.

\bibitem{leray}
J.~Leray, \emph{Sur le mouvement d'un liquide visqueux emplissant l'espace},
Acta Math.\ \textbf{63} (1934), 193--248.

\bibitem{hopf}
E.~Hopf, \emph{\"Uber die Anfangswertaufgabe f\"ur die hydrodynamischen
Grundgleichungen}, Math.\ Nachr.\ \textbf{4} (1951), 213--231.

\bibitem{kato}
T.~Kato, \emph{Strong $L^p$-solutions of the Navier--Stokes equation in
$\R^m$, with applications to weak solutions}, Math.\ Z.\ \textbf{187} (1984),
471--480.

\bibitem{bkm}
J.~T.~Beale, T.~Kato and A.~Majda, \emph{Remarks on the breakdown of smooth
solutions for the 3-D Euler equations}, Comm.\ Math.\ Phys.\ \textbf{94}
(1984), 61--66.

\bibitem{ckn}
L.~Caffarelli, R.~Kohn and L.~Nirenberg, \emph{Partial regularity of suitable
weak solutions of the Navier--Stokes equations}, Comm.\ Pure Appl.\ Math.\
\textbf{35} (1982), 771--831.

\bibitem{fefferman}
C.~L.~Fefferman, \emph{Existence and smoothness of the Navier--Stokes
equation}, in: The Millennium Prize Problems, Clay Math.\ Inst., 2006,
57--67.

\end{thebibliography}

\end{document}

% ledger: placed=28 partial=1 unplaced=44 total=73
